\documentclass[11pt,a4paper,parskip=half,toc=sectionentrywithdots]{scrartcl}
\usepackage[ngerman]{babel}
\usepackage{fontspec}
\usepackage{unicode-math}
\setmainfont{DejaVu Serif}[Scale=0.92]
\setsansfont{DejaVu Sans}[Scale=0.92]
\setmonofont{DejaVu Sans Mono}[Scale=0.80]
\setmathfont{TeX Gyre DejaVu Math}[Scale=0.92]
\usepackage[a4paper,margin=2.3cm]{geometry}
\usepackage{microtype}
\usepackage{xcolor}
\definecolor{linkblue}{HTML}{0B5CAD}
\usepackage{graphicx}
\setkeys{Gin}{width=0.8\linewidth,keepaspectratio}
\usepackage{float}
\floatplacement{figure}{H}
\usepackage{longtable,booktabs,array,calc}
\usepackage{fvextra}
\DefineVerbatimEnvironment{verbatim}{Verbatim}{breaklines,breakanywhere,fontsize=\footnotesize,frame=leftline,framerule=0.6pt,rulecolor=\color{black!30},xleftmargin=4pt}
\usepackage{caption}
\captionsetup{font=small,labelfont=bf}
\usepackage[hidelinks=false,colorlinks=true,linkcolor=linkblue,urlcolor=linkblue,citecolor=linkblue,pdfusetitle,unicode]{hyperref}
\usepackage{xurl}
\usepackage{scrlayer-scrpage}
\pagestyle{scrheadings}
\clearpairofpagestyles
\ohead{\small\headmark}
\cfoot{\small\pagemark}
\automark{section}
\setkomafont{disposition}{\normalfont\bfseries\sffamily}
\setcounter{secnumdepth}{2}
\usepackage{newunicodechar}
\newunicodechar{ℓ}{\ensuremath{\ell}}
\newunicodechar{≳}{\ensuremath{\gtrsim}}
\newunicodechar{≲}{\ensuremath{\lesssim}}
\newunicodechar{≍}{\ensuremath{\asymp}}
\newunicodechar{∖}{\ensuremath{\setminus}}
\providecommand{\tightlist}{\setlength{\itemsep}{0pt}\setlength{\parskip}{0pt}}
\providecommand{\real}[1]{#1}
% tables: small font, a bit less padding
\AtBeginEnvironment{longtable}{\footnotesize\setlength{\tabcolsep}{4pt}\renewcommand{\arraystretch}{1.15}}
\setlength{\LTpre}{6pt}\setlength{\LTpost}{6pt}
\setlength{\emergencystretch}{3em}
\sloppy

\title{Neue Anwendungsfelder für die Graph Zeta Library}
\subtitle{Recherchebericht mit numerischer Nachprüfung der Pilotrechnungen}
\author{Projekt \texttt{GapYas}\thanks{Recherche, Auswertung und Rechnungen erstellt mit Claude Code. Alle Rechnungen sind im Repository reproduzierbar (\texttt{research/verification/}).}}
\date{25.~September 2026}

\begin{document}
\maketitle
\begin{abstract}
\noindent
Dieser Bericht untersucht, wo die im September 2026 veröffentlichte \emph{Graph Zeta Library} (GZL) über ihren ursprünglichen Zweck hinaus neue Ergebnisse liefern kann. Grundlage sind eine Recherche auf arXiv, vor allem zu Arbeiten aus 2025 und 2026, und eigene Rechnungen. Die stärkste Idee ist ein Atlas der Gitter, die Mehrkörper-Potenzgesetzenergien minimieren. Die zugehörigen Pilotrechnungen wurden mit einer von GZL unabhängigen Referenzimplementierung nachgeprüft. Für die Dreikörperenergie $T_\nu$ ergibt sich in 2D ein Phasendiagramm mit drei Phasen: hexagonal $\to$ Quadrat bei $\nu^*=3{,}9183649026$ (erste Ordnung) und Quadrat $\to$ Rechteck bei $\nu_2\approx 8{,}606$ (kontinuierlich). In 3D ist BCC bei allen untersuchten Exponenten das beste gefundene Gitter, während FCC ab $\nu\approx 3{,}75$ instabil wird. Planare modulare Graphfunktionen der Stringtheorie lassen sich als Graph-Zetas des dualen Graphen schreiben; zwei bekannte Identitäten wurden damit auf $10^{-14}$ bis $10^{-15}$ bestätigt.
Eine gezielte Literaturprüfung fand für den 2D-Phasenverlauf, die Minima der zugehörigen modularen Graphfunktionen $C_{s,s,s}$ und das globale 3D-Ergebnis keine Vorarbeit. Die engste Vorarbeit (arXiv:2504.07338) behandelt nur $\nu=3$ und den Bain-Pfad; ihre veröffentlichten Werte stimmen mit unserer Rechnung überein.

\end{abstract}
\tableofcontents
\clearpage
\begin{quote}
\textbf{Hinweis zur Einordnung (Stand 26.09.2026).} Dies ist der breite Recherchebericht, der die Ideen gesammelt hat. Die weitergeführte Forschung zu Idee 1 steht in einem eigenen, fokussierten Dokument (\texttt{research/reports/Dreikörper-Gitterenergie\ –\ Ergebnisse.pdf}) und in einer englischen Note (\texttt{research/note/}). Nach einer externen Kritik wurden hier Überverkaufungen korrigiert. Die 19 Ideen sind Rohmaterial und nicht gleichwertig geprüft.
\end{quote}

\textbf{Kurzfazit (korrigiert).} Die tragfähigste Idee ist Idee 1: die Frage, welches Gitter eine \textbf{reine Dreikörper-Potenzgesetzenergie} T\_ν = Σ'\textbar x\textbar{}\textsuperscript{\{−ν\}\textbar y\textbar{}}\{−ν\}\textbar x−y\textbar{}\textsuperscript{−ν} bei fester Dichte minimiert. Das ist eine \textbf{Modellenergie}, kein Material: Echte Dreikörperkräfte (ATM) haben einen Winkelfaktor und treten immer zusammen mit Paarkräften auf. Die Aussage „BCC schlägt FCC`` gilt daher nur für dieses Modell und steht \textbf{nicht} im Widerspruch zu Sätzen über Paarenergien. Numerisch, und nachgeprüft (Abschnitt 7), durchläuft das 2D-Optimum hexagonal → Quadrat (ν* = 3,9183649026) → Rechteck (ν₂ ≈ 8,606). Im Limes ν → ∞ ist inzwischen \textbf{bewiesen}, dass das Optimum gegen das Rechteck mit Seitenverhältnis √((√17 − 1)/2) läuft. Dass Hexagon und Quadrat kritische Punkte sind, ist dagegen geschenkt (Symmetrie-Fixpunkte). Dass die 2D-Summe eine modulare Graphfunktion C\textsubscript{s,s,s} ist, ist \textbf{per Konstruktion bekannt} und kein Befund; neu ist nur die Frage nach ihrem Minimum. Die Neuheit ist durch Literaturprüfungen gestützt, aber nicht garantiert. Das \textbf{Scoop-Risiko ist hoch}, weil die Gruppe Buchheit/Schwerdtfeger genau diese Funktion definiert hat und Arbeit zur Mehrkörper-Stabilität von Kristallgittern ankündigt. Die übrigen Ideen (Reihenentwicklungen, Quantensimulatoren u. a.) sind unverändert Vorschläge.

Kennzeichnung im ganzen Bericht: \textbf{{[}belegt{]}} heißt, es steht in einer zitierten Quelle. \textbf{{[}abgeleitet{]}} ist eine Schlussfolgerung aus belegten Fakten. \textbf{{[}spekulativ{]}} muss erst geprüft werden. \textbf{{[}Pilot{]}} markiert eigene, \textbf{nicht begutachtete} Rechnungen eines Recherche-Agenten mit GZL 1.0.0 und epsteinlib 0.6.2. \textbf{{[}nachgeprüft{]}} heißt: mit zwei unabhängigen Verfahren übereinstimmend berechnet (Abschnitt 7, Skripte in \texttt{research/verification/}). Die zugehörigen Skripte liegen unter \texttt{research/pilots/} (\texttt{mgf\_test.py}, \texttt{mgf\_test2.py}, \texttt{opt\_test.py}, \texttt{opt2.py}, \texttt{opt3.py}). \texttt{mgf\_test2.py} lädt \texttt{mgf\_test.py} über einen relativen Pfad und muss deshalb aus diesem Verzeichnis gestartet werden.

\begin{center}\rule{0.5\linewidth}{0.5pt}\end{center}

\section{Was GZL wirklich kann: eine Softcore-Gittersumme mit genau einem Impuls}

\subsection{Fünf Sätze zur Klasse der Summen}

\textbf{(1)} GZL berechnet für einen endlichen Multigraphen G mit zwei Terminals s, t die Summe ζ\textsubscript{G}(k) = Σ\_\{x\_v ∈ Λ, v ≠ p\} e\^{}\{−2πi(x\_t − x\textsubscript{s})·k\} Π\textsubscript{e} K\_e(x\textsubscript{e+} − x\textsubscript{e−}). Dabei ist ein Knoten p gepinnt, Λ = A·ℤ\textsuperscript{d} ein Bravais-Gitter mit d = 1, 2, 3, und das Ergebnis ist die Gitter-Fourier-Transformierte eines Zweipunktkorrelators \textbf{{[}belegt{]}} (\href{https://github.com/graph-zeta/gzl}{README}, \href{https://arxiv.org/abs/2609.18918}{arXiv:2609.18918}). \textbf{(2)} Die Summe ist \textbf{softcore}: Nur die beiden Endpunkte derselben Kante werden durch K\_ν(0) = 0 getrennt. Nicht benachbarte Knoten dürfen auf demselben Platz liegen. Die physikalische Hardcore-Bedingung wird in der Physik-Arbeit durch eine kombinatorische „softcore mapping'' in die Korpus-Vorfaktoren eingerechnet, die \textbf{nicht als Code in GZL 1.0.0 enthalten} ist \textbf{{[}belegt{]}} (\href{https://arxiv.org/html/2609.18761}{arXiv:2609.18761}, \href{https://github.com/graph-zeta/gzl}{README}). \textbf{(3)} Die Kerne haben die Form K = a(x) + Σ\textsubscript{j} b\textsubscript{j} \textbar x\textbar{}\textsuperscript{−ν\_j}. Der Anteil a ist reell, gerade, kompakt getragen und darf anisotrop sein, der Potenzgesetzanteil ist isotrop. Jede Kante darf einen eigenen Kern haben, und ν = ∞ liefert den Nächste-Nachbarn-Indikator mit exakten ganzzahligen Einbettungszahlen \textbf{{[}belegt{]}} (\href{https://github.com/graph-zeta/gzl/blob/main/DOCUMENTATION.md}{DOCUMENTATION.md}). \textbf{(4)} Die Rechnung faktorisiert in Blöcke. Brücken sind geschlossene Epstein-Zetas. Einfache Zyklen bei k = 0 laufen über ein Brillouin-Zonen-Integral von Epstein-Produkten. Serien-parallele Blöcke laufen über eine semi-analytische Algebra mit Fehler \textasciitilde{} n\textsuperscript{−(d+σ+2)} und Kosten linear in \textbar V\textbar. Blöcke mit höherer Baumweite laufen über ein Torus-Tensornetz mit Kosten O(N\textsuperscript{tw}), N = n\textsuperscript{d}.~Das ganze k-Gitter kostet per FFT so viel wie ein einzelner Punkt \textbf{{[}belegt{]}} (\href{https://arxiv.org/html/2609.18918}{arXiv:2609.18918}). \textbf{(5)} GZL berechnet \textbf{nicht}: nicht-Bravais-Gitter, anisotrope Potenzgesetzschwänze, komplexe oder ungerade Kerne, oszillierende Schwänze (RKKY) oder Yukawa-Schwänze außer als abgeschnittene Tabelle, Summen mit zwei oder mehr freien Terminals, nicht-Brücken-Blöcke mit ν ≤ d und Ergebnisse jenseits doppelter Genauigkeit \textbf{{[}belegt{]}} (\href{https://github.com/graph-zeta/gzl/blob/main/DOCUMENTATION.md}{DOCUMENTATION.md}).

\subsection{Grenzen und Korrekturen am Auftrag (Abschnitt 3)}

\begin{longtable}[]{@{}
  >{\raggedright\arraybackslash}p{(\columnwidth - 6\tabcolsep) * \real{0.1408}}
  >{\raggedright\arraybackslash}p{(\columnwidth - 6\tabcolsep) * \real{0.2959}}
  >{\raggedright\arraybackslash}p{(\columnwidth - 6\tabcolsep) * \real{0.3602}}
  >{\raggedright\arraybackslash}p{(\columnwidth - 6\tabcolsep) * \real{0.1731}}@{}}
\toprule\noalign{}
\begin{minipage}[b]{\linewidth}\raggedright
Punkt
\end{minipage} & \begin{minipage}[b]{\linewidth}\raggedright
Aussage im Auftrag
\end{minipage} & \begin{minipage}[b]{\linewidth}\raggedright
Tatsächlicher Stand v1.0.0
\end{minipage} & \begin{minipage}[b]{\linewidth}\raggedright
Status
\end{minipage} \\
\midrule\noalign{}
\endhead
\bottomrule\noalign{}
\endlastfoot
Summationsbereich & „Platzierungen der Knoten auf Λ'' & Softcore: Nur Kantenendpunkte sind getrennt. Hardcore→Softcore-Abbildung ist nicht mitgeliefert. & \textbf{Korrektur/Präzisierung} \\
Exponenten nahe d & „präzise auch für ν nahe d'' & Stimmt nur für ν \textgreater{} d.~Brücken werden meromorph fortgesetzt. \textbf{Jeder andere Block mit ν ≤ d wirft \texttt{UnsupportedLatticeSumError}.} Benchmarks reichen bis σ = ν − d = 0,5. & \textbf{Einschränkung} \\
„analytisch'' & Brücken und Zyklen analytisch & Brücken sind geschlossen. Zyklen nur bei k = 0 über \texttt{zeta\_circle}, und nur für d ≤ 3. SP-Blöcke sind semi-analytisch mit algebraischem Fehler. & Präzisierung \\
Kosten tw ≥ 3 & „polynomial in der Gittergröße'' & O(N\textsuperscript{tw}) Operationen und O(N\textsuperscript{tw−1}) Speicher. \textbf{In d = 3 ist praktisch n\textsubscript{points} ≤ 16.} Bei n = 8 liegt der relative Fehler hoher Ordnungen um 10⁻³. & Präzisierung \\
\texttt{slab\_zeta} & klingt nach Film-/Slab-Geometrie & \textbf{Ist keine Slab-Geometrie}, sondern eine speichersparende Engine für den Vakuumwert eines dichten Blocks & \textbf{Korrektur} \\
Kerne & reell, gerade & Stimmt. Die Mathe-Arbeit erlaubt komplexe ℓ¹-Kerne, \textbf{die Bibliothek ist enger als die Theorie}. & Präzisierung \\
Anisotropie & Roadmap & Stimmt. KTmSe₂ in 3D wurde mit einer experimentellen Erweiterung nur bis Ordnung 5 gerechnet. Die anisotrope Epstein-Zeta ist in EpsteinLib bereits vorhanden (\href{https://arxiv.org/abs/2609.28282}{arXiv:2609.28282}). & bestätigt \\
Dimension & d = 1, 2, 3 & Offiziell ja. d ≥ 4 läuft teilweise (Brücken exakt), ist aber unkalibriert und nicht publikationsfähig. & Ergänzung \\
Gittermatrix & jede Matrix A & A muss quadratisch und regulär sein, also keine eingebetteten niedrigdimensionalen Gitter & Präzisierung \\
Impuls & e\textsuperscript{ik·(\ldots)} & Genau \textbf{ein} externer Impuls, in fraktionalen reziproken Koordinaten. Zwei oder mehr freie Terminals werden abgelehnt, 2qp-Sektoren gibt es also nicht. & \textbf{Einschränkung} \\
Korpora & TFIM 0qp bis Ordnung 13, 1qp bis Ordnung 11 & Bestätigt: 8.403 bzw. 22.677 Graphen. Die Korpora sind dimensionsunabhängig. & bestätigt \\
\end{longtable}

Die Leistungsbehauptungen halten der Prüfung stand \textbf{{[}belegt{]}}. 114 Reihen mit 998 Koeffizienten stimmen mit publizierten MC-Daten überein: \textbf{maximal 3,1 σ\textsubscript{MC}, Median 0,12 σ\textsubscript{MC}}. Der Satz wurde in \textbf{5 Minuten auf 8 M1-Kernen} gerechnet, gegenüber etwa \textbf{3×10⁴ Kernstunden pro MC-Reihe und Impuls}. Eine 3D-Dispersion der Ordnung 10 auf 16³ Impulspunkten dauert rund 10 Minuten auf einem Kern (\href{https://arxiv.org/html/2609.18918}{arXiv:2609.18918}). Die Physik-Arbeit liefert 0qp- und 1qp-Reihen in 3D bis Ordnung 13 bzw. 11 in etwa 200 s (\href{https://arxiv.org/html/2609.18761}{arXiv:2609.18761}).

Stand der Verbreitung am 25.09.2026: 1 Stern, 0 Forks, 0 Issues, keine Nutzung außerhalb der Gruppe Buchheit/Schmidt. \textbf{Das Feld ist also offen.} Die Gruppe Sbierski/Lesanovsky ist als einzige externe Nutzerin der Vorläufer-Werkzeuge nachweisbar und wäre die wahrscheinlichste erste Konkurrenz (\href{https://arxiv.org/abs/2604.22743}{arXiv:2604.22743}).

\subsection{Scoop-Landkarte: was die GZL-Autoren selbst angekündigt haben}

Die Autoren beanspruchen öffentlich folgende Richtungen \textbf{{[}belegt{]}} (\href{https://arxiv.org/html/2609.18918}{arXiv:2609.18918 §10}, \href{https://arxiv.org/html/2609.18761}{arXiv:2609.18761 §VIII}):

\begin{itemize}
\tightlist
\item
  anisotrope Dipolkerne, konkret LiHoF₄, Fe₈, Mn₁₂ und RE(OH)₃;
\item
  mehratomige Gitter und frustrierte Systeme;
\item
  Heisenberg- und weitere Modelle;
\item
  Spektralgewichte und höhere Quasiteilchensektoren;
\item
  endliche Temperatur, ausdrücklich mit Burkard, Schneider und Sbierski;
\item
  korrelierte Unordnung;
\item
  Rydberg-vdW-Arrays im Vergleich mit dynamischen Strukturfaktoren;
\item
  eine analytische Singularitätsbehandlung für die Tensornetz-Blöcke.
\end{itemize}

Wer dort ohne Kooperation arbeitet, läuft Gefahr, überholt zu werden. \textbf{Nicht genannt} sind:

\begin{itemize}
\tightlist
\item
  klassische Hochtemperatur- und Perkolationsreihen;
\item
  Gitteroptimierung;
\item
  modulare Graphfunktionen;
\item
  Gitter-Störungstheorie.
\end{itemize}

Im Umfeld der Autoren liegen allerdings ATM-Energien entlang des Bain-Pfads (\href{https://arxiv.org/abs/2504.07338}{arXiv:2504.07338}, \href{https://arxiv.org/abs/2504.11989}{arXiv:2504.11989}).

\begin{center}\rule{0.5\linewidth}{0.5pt}\end{center}

\section{Top-5-Projektideen: Zwei Transfer-Ideen führen, zwei Physik-Ideen sind experimentnah}

\begin{longtable}[]{@{}
  >{\raggedright\arraybackslash}p{(\columnwidth - 20\tabcolsep) * \real{0.0490}}
  >{\raggedright\arraybackslash}p{(\columnwidth - 20\tabcolsep) * \real{0.3338}}
  >{\raggedright\arraybackslash}p{(\columnwidth - 20\tabcolsep) * \real{0.0490}}
  >{\raggedright\arraybackslash}p{(\columnwidth - 20\tabcolsep) * \real{0.0642}}
  >{\raggedright\arraybackslash}p{(\columnwidth - 20\tabcolsep) * \real{0.0642}}
  >{\raggedright\arraybackslash}p{(\columnwidth - 20\tabcolsep) * \real{0.0490}}
  >{\raggedright\arraybackslash}p{(\columnwidth - 20\tabcolsep) * \real{0.0921}}
  >{\raggedright\arraybackslash}p{(\columnwidth - 20\tabcolsep) * \real{0.0490}}
  >{\raggedright\arraybackslash}p{(\columnwidth - 20\tabcolsep) * \real{0.0785}}
  >{\raggedright\arraybackslash}p{(\columnwidth - 20\tabcolsep) * \real{0.0490}}
  >{\raggedright\arraybackslash}p{(\columnwidth - 20\tabcolsep) * \real{0.0921}}@{}}
\toprule\noalign{}
\begin{minipage}[b]{\linewidth}\raggedright
\#
\end{minipage} & \begin{minipage}[b]{\linewidth}\raggedright
Idee
\end{minipage} & \begin{minipage}[b]{\linewidth}\raggedright
Spur
\end{minipage} & \begin{minipage}[b]{\linewidth}\raggedright
Neuheit
\end{minipage} & \begin{minipage}[b]{\linewidth}\raggedright
Passung
\end{minipage} & \begin{minipage}[b]{\linewidth}\raggedright
Hebel
\end{minipage} & \begin{minipage}[b]{\linewidth}\raggedright
Machbarkeit
\end{minipage} & \begin{minipage}[b]{\linewidth}\raggedright
Wow
\end{minipage} & \begin{minipage}[b]{\linewidth}\raggedright
Anschluss
\end{minipage} & \begin{minipage}[b]{\linewidth}\raggedright
Σ
\end{minipage} & \begin{minipage}[b]{\linewidth}\raggedright
Scoop
\end{minipage} \\
\midrule\noalign{}
\endhead
\bottomrule\noalign{}
\endlastfoot
1 & Optimale Gitter für Mehrkörper-Potenzgesetzenergien & B & 5 & 5 & 4 & 5 & 5 & 4 & \textbf{28} & \textbf{hoch} \\
2 & Klassische LR-Hochtemperatur- und Perkolationsreihen gegen den Sak-Streit & B & 4 & 3 & 4 & 2 & 3 & 5 & \textbf{21} & mittel \\
3 & Modulare Graphfunktionen als duale Graph-Zetas & B & 5 & 3 & 2 & 4 & 4 & 3 & \textbf{21} & gering \\
4 & Simulator-Material-Zwillinge: Rydberg-vdW und Tm-Dreiecksmagnete & A & 3 & 4 & 4 & 4 & 5 & 4 & \textbf{24} & \textbf{hoch} \\
5 & Quanten-Sak-Test und Log-Korrekturen im LRTFIM (d = 1--3) & A & 3 & 5 & 3 & 4 & 3 & 4 & \textbf{22} & mittel--hoch \\
\end{longtable}

Die Reihenfolge gewichtet die Neuheit stärker als die bloße Punktsumme, weil der Auftrag „etwas völlig Neues'' verlangt.

\subsection{Idee 1: Optimale Gitter für Mehrkörper-Potenzgesetzenergien („Mehrkörper-Sarnak--Strömbergsson'')}

\textbf{Pitch.} GZL als erstes Werkzeug, das die Frage „Welches Gitter minimiert eine Drei- oder Vierkörper-Potenzgesetzenergie?'' über den \emph{gesamten} Modulraum der 2D- und 3D-Gitter in Sekunden pro Punkt beantwortet. Die Pilotrechnungen deuten auf eine andere Antwort als im Zweikörperfall.

\textbf{Das Problem im Zielfeld.} Die Optimierung von Gitterenergien ist in der Mathematik gerade ein heißes Thema. Luo und Wei beanspruchen einen Beweis der Sarnak--Strömbergsson-Vermutung: FCC minimiert die 3D-Epstein-Zeta für alle s \textgreater{} 3/2. Der Beweis ist zehn Tage alt und noch nicht begutachtet \textbf{{[}belegt{]}} (\href{https://arxiv.org/abs/2609.17356}{arXiv:2609.17356}). Dazu kommen Klassifikationen von Minimierern für Verhältnisse von Theta- und Epstein-Funktionen (\href{https://arxiv.org/abs/2605.07580}{arXiv:2605.07580}) und überraschende Phasenübergänge der Minimierer bei p-Normen (\href{https://arxiv.org/abs/2407.20762}{arXiv:2407.20762}). Die ganze Linie behandelt aber \textbf{Zweikörper-Energien}, also Summen über eine Kante. In der Kristallchemie ist dagegen belegt, dass Dreikörperterme (ATM) die Phasenstabilität verschieben. Starke repulsive ATM-Beiträge können BCC gegenüber FCC begünstigen, doch „the bcc phase remains susceptible to further cuboidal distortions''. Untersucht wurde dort nur der Bain-Pfad, also eine Ein-Parameter-Kurve \textbf{{[}belegt{]}} (\href{https://arxiv.org/abs/2504.07338}{arXiv:2504.07338}, \href{https://arxiv.org/abs/2504.11989}{arXiv:2504.11989}). Eine Websuche nach Dreikörper-Gitteroptimierung über alle 2D- und 3D-Gitter fand keine Arbeit aus 2025 \textbf{{[}belegt als Negativbefund{]}}. Der Grund für die Lücke ist schlicht Rechenaufwand: Eine Dreikörpersumme ist eine Doppelsumme über ein unendliches Gitter mit langsamem Abfall nahe ν = d.

\textbf{Abbildung auf GZL.} Die Dreikörper-Riesz-Energie E₃(Λ) = Σ'\textsubscript{x,y} \textbar x\textbar{}\textsuperscript{\{−ν\}\textbar y\textbar{}}\{−ν\}\textbar x−y\textbar{}\textsuperscript{−ν} ist \textbf{exakt der Dreiecksgraph (3-Zyklus)}. Vierkörperenergien sind der 4-Zyklus oder K₄ (tw = 3, Tensornetz). Das Gitter A ist der freie Parameter. In 2D ist der Modulraum die Fundamentaldomäne (2 Parameter), in 3D sind es 5 Parameter (Gitter modulo Rotation und Skala) \textbf{{[}abgeleitet{]}}.

\begin{verbatim}
import numpy as np, gzl, epsteinlib
def A2(tau):                        # 2D-Gitter mit Einheitsvolumen, tau in der Fundamentaldomäne
    A = np.array([[1, tau.real], [0, tau.imag]]); return A/np.sqrt(tau.imag)
E2 = lambda A, nu: epsteinlib.epstein_zeta(nu, A, np.zeros(2), np.zeros(2)).real  # Zweikörper
E3 = lambda A, nu: gzl.zeta_circle([nu, nu, nu], A)                               # Dreieck
K4 = np.array([[0,1],[0,2],[0,3],[1,2],[1,3],[2,3]])
E4 = lambda A, nu: gzl.evaluate_graph(K4, np.full(6, nu), A, n_points=48)         # Vierkörper, tw=3
# Observable: argmin_A [c2*E2 + c3*E3 + c4*E4] über den Modulraum -> Phasendiagramm in (nu, c3/c2)
\end{verbatim}

\textbf{Pilotergebnisse {[}Pilot, inzwischen nachgeprüft und ergänzt, siehe Abschnitt 7{]}} (\texttt{opt\_test.py}, \texttt{opt2.py}, \texttt{opt3.py}; GZL \texttt{zeta\_circle}, Einheitsvolumen, gleiche ν auf allen drei Kanten; \href{https://pypi.org/project/gzl/}{gzl auf PyPI}):

In 2D ist das Hexagonalgitter für ν ≲ 3,918 global optimal, darüber das Quadratgitter. \textbf{Korrektur aus der Nachprüfung:} Das Quadrat bleibt nur bis ν₂ ≈ 8,606 optimal, danach gewinnt ein Rechteckgitter (Abschnitt 7.2).

\begin{longtable}[]{@{}llll@{}}
\toprule\noalign{}
ν & Globales Minimum & Hexagonal & Quadrat \\
\midrule\noalign{}
\endhead
\bottomrule\noalign{}
\endlastfoot
2,5 & hex & 20,6364817949 & 21,0272603756 \\
3,0 & hex & 13,3917941102 & 13,6528937154 \\
3,5 & hex & 9,7723761021 & 9,8932825030 \\
3,75 & hex & 8,5685550844 & 8,6172422133 \\
4,0 & Quadrat & 7,6061707526 & 7,5826699011 \\
6,0 & Quadrat & 3,7457712638 & 3,2499743124 \\
\end{longtable}

Der Schnittpunkt liegt bei \textbf{ν* ≈ 3,91836} (Brent-Nullstelle). Bei ν* steigt die Energie entlang des Randbogens \textbar τ\textbar{} = 1 von 7,898512 (hex, 60°) auf etwa \textbf{7,9096 bei \textasciitilde72°} und fällt wieder auf 7,898512 (Quadrat, 90°). Es gibt also eine Barriere, und der Übergang ist \textbf{erster Ordnung}. Ein direkter Brute-Force-Vergleich auf ℤ² bei ν = 6 ergibt 3,2499743122 gegenüber 3,2499743124 aus GZL. Die Erklärung (eigene Abschätzung): Bei großem ν dominieren die kleinsten Dreiecke, und repulsive Dreikörperterme bestrafen die zwölf gleichseitigen Nächste-Nachbarn-Dreiecke des Hexagonalgitters.

In 3D (nur FCC, BCC und SC verglichen, \textasciitilde28 s pro Satz) liegt \textbf{BCC bei allen getesteten ν vor FCC}:

\begin{longtable}[]{@{}llll@{}}
\toprule\noalign{}
ν & FCC & BCC & SC \\
\midrule\noalign{}
\endhead
\bottomrule\noalign{}
\endlastfoot
3,5 & 38,1993 & \textbf{38,0744} & 41,9368 \\
4,5 & 17,9624 & \textbf{17,7406} & 21,0981 \\
6,0 & 8,0434 & \textbf{7,7102} & 10,2962 \\
9,0 & 2,3133 & \textbf{1,9785} & 3,2723 \\
\end{longtable}

\textbf{Was neu herauskäme.} Die erste Karte der Gitteroptima für reine Mehrkörperenergien über den gesamten Modulraum, mit drei Bausteinen:

\begin{itemize}
\tightlist
\item
  Phasendiagrammen in ν und im Kopplungsverhältnis Zwei- zu Dreikörper;
\item
  einer belastbar formulierten Vermutung „BCC minimiert die Dreiecks-Zeta in 3D'' als Mehrkörper-Gegenstück zu Sarnak--Strömbergsson;
\item
  dem ersten dokumentierten Hexagonal-Quadrat-Übergang erster Ordnung für eine reine Dreikörperenergie.
\end{itemize}

\textbf{{[}abgeleitet{]}} Für ganzzahlige Exponenten gibt es beweisbare Ankerpunkte. Die 2D-Energien mit Exponenten (2,2,2) und (4,4,2) reduzieren sich über die Identitäten C₁,₁,₁ = E₃ + ζ(3) und C₂,₂,₁ = (2/5)E₅ + ζ(5)/30 auf Epstein-Zetas. Diese werden bekanntlich vom Hexagonalgitter minimiert (Hintergrundwissen, in den Notizen nicht belegt). Das passt zur Pilotaussage „hex für kleine ν'' und eröffnet einen Beweisweg über modulare Graphfunktionen (Idee 3).

\textbf{Erster Meilenstein (2--4 Wochen).} Die Piloten mit Konvergenzstudie reproduzieren: n\textsubscript{points} variieren, \texttt{zeta\_circle} gegen \texttt{evaluate\_graph} prüfen, Brute-Force-Vergleiche bei mehreren ν. Die 2D-Suche auf große Im τ ausdehnen und ν* samt Barriere bestätigen.

\textbf{Ausbaupfad.}

\begin{itemize}
\tightlist
\item
  \textbf{Monat 2--3:} globale 3D-Suche im 5-dimensionalen Modulraum, zum Beispiel über reduzierte Gram-Matrizen, auf einem ν-Gitter von 3,2 bis 12. Ergänzend 4-Zyklus und K₄.
\item
  \textbf{Monat 4--6:} gemischte Energien aus Zweikörper-Riesz bzw. Lennard-Jones und Dreieck mit Phasendiagramm. Damit direkt die offene Frage aus \href{https://arxiv.org/abs/2504.07338}{arXiv:2504.07338} angehen: Ist BCC gegen \emph{nicht-kuboidale} Verzerrungen stabil?
\item
  \textbf{Monat 6--12:} ATM über Zerlegung in Dreiecks-Zetas. Dafür ist eine \textbf{GZL-Erweiterung nötig}: meromorphe Fortsetzung für Dreiecksblöcke mit ν ≤ d, siehe Risiken. Außerdem Beweisversuche für ganzzahlige s in 2D über MGF-Laplace-Gleichungen.
\end{itemize}

\textbf{Vorwissen und Einstieg.} Gittergeometrie (Fundamentaldomäne, Reduktionstheorie), globale Optimierung, Grundlagen der Epstein-Zeta. Einstieg über \href{https://arxiv.org/abs/2609.17356}{arXiv:2609.17356}, \href{https://arxiv.org/abs/2312.01395}{arXiv:2312.01395} (Strukturübergänge auf Rechteckgittern) und \href{https://arxiv.org/abs/2504.11989}{arXiv:2504.11989}.

\textbf{Risiken.}

\begin{enumerate}
\def\labelenumi{\arabic{enumi}.}
\tightlist
\item
  Neuheit ist nur über das Fehlen von Treffern gesichert. Die Bétermin-Gruppe (offene Fragen, zitiert in \href{https://arxiv.org/abs/2411.17199}{arXiv:2411.17199}) und der Schwerdtfeger-Buchheit-Kreis könnten Ähnliches in Arbeit haben.
\item
  Die 2D-Suche deckte nur Im τ ≤ 2--3 mit sechs Startpunkten ab. In 3D wurden nur drei Gitter verglichen.
\item
  ATM enthält pro Kante Exponenten ν ∈ \{5, 3, 1, −1\}, also ν ≤ d.~\textbf{Das lehnt GZL 1.0.0 für Nicht-Brücken ab.} Die Fortsetzung existiert für Zyklen im Buchheit-Busse-Verfahren, ist in GZL aber nicht freigelegt.
\item
  Die Zahlen sind inzwischen unabhängig nachgerechnet (Abschnitt 7), aber nicht begutachtet. Die 3D-Suche mit 16 Starts ist kein Beweis des globalen Minimums.
\end{enumerate}

\textbf{Scoop-Risiko: hoch (korrigiert).} Die GZL-Roadmap nennt Gitteroptimierung nicht, aber der Autorenkreis hat genau diese Funktion definiert, arbeitet an ATM-Energien auf dem Bain-Pfad und kündigt Arbeit zur Mehrkörper-Stabilität an. Eine frühe Kontaktaufnahme mit Buchheit und Schwerdtfeger könnte aus dem Risiko eine Kooperation machen.

\textbf{Schlüsselpapers.}

\begin{longtable}[]{@{}
  >{\raggedright\arraybackslash}p{(\columnwidth - 6\tabcolsep) * \real{0.1032}}
  >{\raggedright\arraybackslash}p{(\columnwidth - 6\tabcolsep) * \real{0.4037}}
  >{\raggedright\arraybackslash}p{(\columnwidth - 6\tabcolsep) * \real{0.0593}}
  >{\raggedright\arraybackslash}p{(\columnwidth - 6\tabcolsep) * \real{0.4037}}@{}}
\toprule\noalign{}
\begin{minipage}[b]{\linewidth}\raggedright
arXiv
\end{minipage} & \begin{minipage}[b]{\linewidth}\raggedright
Titel / Inhalt (laut Notizen)
\end{minipage} & \begin{minipage}[b]{\linewidth}\raggedright
Jahr
\end{minipage} & \begin{minipage}[b]{\linewidth}\raggedright
Relevanz
\end{minipage} \\
\midrule\noalign{}
\endhead
\bottomrule\noalign{}
\endlastfoot
\href{https://arxiv.org/abs/2609.17356}{2609.17356} & Luo \& Wei, beanspruchter Beweis der 3D-Sarnak--Strömbergsson-Vermutung & 2026 & Zweikörper-Referenz: FCC optimal, Gegenpol zum BCC-Pilot \\
\href{https://arxiv.org/abs/2504.07338}{2504.07338} & Robles-Navarro et al., LJ + ATM auf kuboidalen (Bain-)Gittern & 2025 & ATM begünstigt BCC. Offene Stabilitätsfrage außerhalb des Bain-Pfads \\
\href{https://arxiv.org/abs/2504.11989}{2504.11989} & Buchheit \& Busse, „Epstein zeta method for many-body lattice sums'' & 2025 & Methodischer Vorläufer, ATM „from weeks to minutes'' \\
\href{https://arxiv.org/abs/2407.20762}{2407.20762} & Bétermin \& Furlanetto, Minimierer für p-Normen & 2024 & Präzedenz für unerwartete Minimierer-Übergänge \\
\href{https://arxiv.org/abs/2312.01395}{2312.01395} & Bétermin, Šamaj, Travěnec, Strukturübergänge auf 2D-Rechteckgittern & 2023 & Theorie von Übergängen erster und zweiter Ordnung, Vorlage für die Analyse \\
\href{https://arxiv.org/abs/2605.07580}{2605.07580} & Luo \& Wei, Minimierer von Verhältnissen und Differenzen von Theta- und Epstein-Funktionen & 2026 & Aktuelle Technik, zentrale Rolle des Hexagonalgitters \\
\end{longtable}

\textbf{Bewertung.}

\begin{longtable}[]{@{}
  >{\raggedright\arraybackslash}p{(\columnwidth - 4\tabcolsep) * \real{0.1851}}
  >{\raggedright\arraybackslash}p{(\columnwidth - 4\tabcolsep) * \real{0.1140}}
  >{\raggedright\arraybackslash}p{(\columnwidth - 4\tabcolsep) * \real{0.6709}}@{}}
\toprule\noalign{}
\begin{minipage}[b]{\linewidth}\raggedright
Kriterium
\end{minipage} & \begin{minipage}[b]{\linewidth}\raggedright
Punkte
\end{minipage} & \begin{minipage}[b]{\linewidth}\raggedright
Begründung
\end{minipage} \\
\midrule\noalign{}
\endhead
\bottomrule\noalign{}
\endlastfoot
Neuheit & 5 & Kein Treffer für Mehrkörper-Optimierung über den ganzen Modulraum. Der Pilot liefert zwei neue qualitative Befunde. \\
Passung & 5 & Läuft mit \texttt{zeta\_circle} und \texttt{evaluate\_graph} in v1.0.0, ν \textgreater{} d.~Nur ATM braucht eine Erweiterung. \\
Hebel & 4 & Sekunden pro Gitterpunkt statt Doppelsummen mit langsamer Konvergenz. Globale Suche wird erst dadurch praktikabel. \\
Machbarkeit & 5 & Die Piloten existieren, der erste Meilenstein ist in Wochen erreichbar. \\
Wow & 5 & Heatmap über die Fundamentaldomäne, springendes Optimum, BCC schlägt FCC. \\
Anschluss & 4 & math.MG, math-ph und Kristallchemie. Mit ATM auch materialrelevant. \\
\end{longtable}

\subsection{Idee 2: Klassische langreichweitige Hochtemperatur- und Perkolationsreihen für den Sak-Streit}

\textbf{Pitch.} Die klassische Schwester der Quanten-Linked-Cluster-Entwicklung, für die GZL gebaut wurde: deterministische Hochtemperaturreihen für LR-Ising, LR-O(n) und LR-Perkolation. Sie liefern β\textsubscript{c}(σ), ρ\textsubscript{c}(σ) und χ(k) über die ganze Brillouin-Zone und geben im Sak-Streit eine Stimme, die nicht auf Monte Carlo beruht.

\textbf{Das Problem im Zielfeld.} Wo die Grenze zwischen lang- und kurzreichweitiger Universalität liegt, ist 2025--2026 offen umstritten \textbf{{[}belegt{]}}:

\begin{itemize}
\tightlist
\item
  Die Deng-Gruppe findet mit MC bis L = 8192 \textbf{σ* = 2} für 2D LR-Ising, XY, Heisenberg und Perkolation, gegen Saks σ* = 2 − η\textsubscript{SR} = 7/4 (\href{https://arxiv.org/abs/2512.04805}{arXiv:2512.04805}, \href{https://arxiv.org/abs/2512.01956}{arXiv:2512.01956}).
\item
  Ein Quanten-RG-Papier beansprucht σ* = 2 auch für Quanten-O(n)-Modelle (\href{https://arxiv.org/abs/2606.22407}{arXiv:2606.22407}).
\item
  In der 2D-LR-Perkolation bis L = 16384 weicht η nahe σ ≈ 3/2 von 2 − σ ab. Dort sind ρ\textsubscript{c}(σ = 1) = 0,307591(4) und ρ\textsubscript{c}(σ = 1/2) = 0,26263(2) bestimmt (\href{https://arxiv.org/abs/2608.20750}{arXiv:2608.20750}).
\item
  Hutchcroft nennt die Berechnung von α\textsubscript{c}(d) für 2 \textless{} d \textless{} 6 „beyond the scope of current techniques'' (\href{https://arxiv.org/abs/2508.18808}{arXiv:2508.18808}).
\end{itemize}

Heute gibt es drei Zugänge: rigorose Schranken ohne Konstanten, MC auf endlichen Tori mit Minimum-Image-Abschneidung und Kontinuums-ε-Entwicklungen (\href{https://arxiv.org/abs/2602.07818}{arXiv:2602.07818}, \href{https://arxiv.org/abs/2608.15120}{arXiv:2608.15120}). \textbf{Eine Reihe mit exakten Gitterkoeffizienten für klassische LR-Modelle wurde für 2025--26 nicht gefunden} \textbf{{[}belegt als Negativbefund{]}}.

\textbf{Abbildung auf GZL {[}abgeleitet{]}.} In der Free-Embedding-Linked-Cluster-Entwicklung ist der Koeffizient von βⁿ in χ(k) oder ln Z eine Summe Σ\textsubscript{G} w\textsubscript{G} · Σ\textsubscript{x} Π\textsubscript{e} J(x\_e)\textsuperscript{m\_e} e\textsuperscript{ik·Δx}, mit J(x) = \textbar x\textbar{}\textsuperscript{−(d+σ)}, J(0) = 0 und Kantenvielfachheit m\textsubscript{e}. Da J\textsuperscript{m} wieder ein Potenzgesetz mit ν = m(d+σ) \textgreater{} d ist, ist das \textbf{exakt eine Graph-Zeta mit GZLs Konvention K(0) = 0}. Die Vertexfaktoren (Kumulanten für Ising, O(n), φ⁴) bringt der Nutzer mit. Für LR-Perkolation liefert die Entwicklung 1 − e\textsuperscript{−βJ} = Σ\textsubscript{m} (−1)\textsuperscript{m+1} β\textsuperscript{m} J\textsuperscript{m}/m! pro Ordnung endliche Summen von Potenzgesetzen, also genau GZLs Kernklasse (\href{https://arxiv.org/abs/2508.18807}{arXiv:2508.18807}). Die Subgraph-Entwicklung verlangt allerdings Hardcore-Einbettungen und damit die Kontraktionsabbildung aus \href{https://arxiv.org/html/2609.18761}{arXiv:2609.18761}, die neu implementiert werden muss.

\begin{verbatim}
# Pseudocode – Korpus und Vertexgewichte sind Eigenbau
for n in range(1, n_max + 1):
    for G, w_G, mult in classical_LR_corpus(model="ising", order=n):   # Graphenenumeration + Kumulanten
        nu_e = [m * (d + sigma) for m in mult]                           # J^m bleibt ein Potenzgesetz
        chi_k[n] += w_G * gzl.evaluate_graph(G.edges, nu_e, A, terminal=(s, t), n_points=N)  # ganze BZ per FFT
# Analyse: Ratio/Padé für beta_c(sigma), gamma(sigma); nichtanalytischer |k|^sigma-Koeffizient von 1/chi(k)
\end{verbatim}

\textbf{Was neu herauskäme.}

\begin{itemize}
\tightlist
\item
  β\textsubscript{c}(σ) und ρ\textsubscript{c}(σ) als glatte Funktionen auf dichten σ-Gittern, bis nahe σ → 0, wo direkte Summation versagt.
\item
  γ(σ) aus Ratio- und Padé-Analyse zum Vergleich mit ε-Entwicklungen und FRG (\href{https://arxiv.org/abs/2510.02458}{arXiv:2510.02458}).
\item
  Die \textbf{Amplitude des \textbar k\textbar\^{}σ-Terms} von χ(k)\textsuperscript{−1}. Für σ \textless{} 2 divergiert das zweite Moment, sodass die Standard-Korrelationslänge der Hochtemperaturreihe gar nicht definiert ist. Das volle k-Gitter aus GZL ist dafür das natürliche Werkzeug \textbf{{[}abgeleitet{]}}.
\item
  Als spekulative Erweiterung: die Edwards-Anderson-Suszeptibilität der 1D-LR-Spinglas-Kette mit J² ∝ r\textsuperscript{−2σ}, für das offene Fenster 3/2 ≤ α ≤ 2 (\href{https://arxiv.org/abs/2604.07130}{arXiv:2604.07130}).
\end{itemize}

\textbf{Erster Meilenstein (6--10 Wochen).}

\begin{enumerate}
\def\labelenumi{\arabic{enumi}.}
\tightlist
\item
  \textbf{Literaturschranke:} Prüfen, ob alte LR-Reihen (Nagle--Bonner 1970, Glumac--Uzelac um 1989; beide in den Notizen \emph{nicht verifiziert}) schon hohe Ordnungen erreichten.
\item
  Eine eigene Graphenenumeration mit Ising-Kumulanten schreiben und χ(k = 0) für 1D bis etwa Ordnung 8 berechnen.
\item
  Niedrige Ordnungen gegen direkte Summen prüfen.
\end{enumerate}

\textbf{Ausbaupfad.}

\begin{itemize}
\tightlist
\item
  \textbf{Monat 3--6:} 2D auf dem Quadrat- und Dreiecksgitter mit dichtem σ-Gitter, Vergleich mit den kritischen Punkten der Deng-Gruppe.
\item
  \textbf{Monat 6--9:} χ(k) über die ganze Brillouin-Zone und die \textbar k\textbar\^{}σ-Amplitude.
\item
  \textbf{Monat 9--15:} LR-Perkolation mit Hardcore-Kontraktionen, Vergleich mit ρ\textsubscript{c} = 0,307591(4). Dabei muss die Normierungskonvention angeglichen werden: Unendliches Gitter und Epstein-Zeta stehen gegen Torus mit Minimum-Image.
\end{itemize}

\textbf{Vorwissen und Einstieg.} Klassische Linked-Cluster-Technik nach Wortis, Englert, Lüscher--Weisz (Lehrbuchstoff, in den Notizen ohne arXiv-Quelle), Graphenenumeration, Reihenanalyse. Einstieg: \href{https://arxiv.org/abs/2403.00421}{arXiv:2403.00421} (MC-Einbettungs-Review, zeigt die Quanten-Pipeline), \href{https://arxiv.org/abs/2512.04805}{arXiv:2512.04805}, \href{https://arxiv.org/html/2608.20750}{arXiv:2608.20750}.

\textbf{Risiken.}

\begin{enumerate}
\def\labelenumi{\arabic{enumi}.}
\tightlist
\item
  Die Graphenzahl bei Ordnung 15--20 wird grob auf 10⁵--10⁷ geschätzt (\emph{nicht verifiziert}). Die Baumweite hoher 1PI-Graphen kann 4--6 erreichen, und das Tensornetz kostet dann O(N\textsuperscript{tw}).
\item
  Nahe σ ≈ 2 machen Log-Korrekturen die Reihenanalyse unzuverlässig. GZL entfernt das Koeffizientenrauschen, nicht den Extrapolationsfehler begrenzter Ordnung.
\item
  Der Hardcore-Code fehlt in GZL.
\item
  Alte Literatur könnte die Neuheit schmälern.
\end{enumerate}

\textbf{Scoop-Risiko: mittel.} Klassische LR-Modelle nennen die Autoren nicht. „Endliche Temperatur'' mit der Gruppe Sbierski steht aber auf ihrer Roadmap, und die Sbierski-Gruppe arbeitet bereits mit Hochtemperaturentwicklungen (\href{https://arxiv.org/abs/2604.22743}{arXiv:2604.22743}).

\textbf{Schlüsselpapers.}

\begin{longtable}[]{@{}
  >{\raggedright\arraybackslash}p{(\columnwidth - 6\tabcolsep) * \real{0.1058}}
  >{\raggedright\arraybackslash}p{(\columnwidth - 6\tabcolsep) * \real{0.4138}}
  >{\raggedright\arraybackslash}p{(\columnwidth - 6\tabcolsep) * \real{0.0608}}
  >{\raggedright\arraybackslash}p{(\columnwidth - 6\tabcolsep) * \real{0.3896}}@{}}
\toprule\noalign{}
\begin{minipage}[b]{\linewidth}\raggedright
arXiv
\end{minipage} & \begin{minipage}[b]{\linewidth}\raggedright
Titel / Inhalt
\end{minipage} & \begin{minipage}[b]{\linewidth}\raggedright
Jahr
\end{minipage} & \begin{minipage}[b]{\linewidth}\raggedright
Relevanz
\end{minipage} \\
\midrule\noalign{}
\endhead
\bottomrule\noalign{}
\endlastfoot
\href{https://arxiv.org/abs/2512.04805}{2512.04805} & Xiao, Liu, Fan, Deng, 2D LR-Ising bis L = 8192, σ* = 2 & 2025 & Hauptgegenstand des Sak-Streits \\
\href{https://arxiv.org/abs/2608.20750}{2608.20750} & Liu, Xiao, Fan, Deng, 2D LR-Perkolation bis L = 16384 & 2026 & Präzise ρ\textsubscript{c-Benchmarks} \\
\href{https://arxiv.org/abs/2606.22407}{2606.22407} & „Perturbative Renormalization and Universality Diagram for Long-Range Quantum Criticality'' & 2026 & Quanten-Gegenstück σ* = 2 \\
\href{https://arxiv.org/abs/2508.18808}{2508.18808} & Hutchcroft, „Critical LRP II'' & 2025 & Offene Frage α\textsubscript{c}(d), nur Aussagen bis auf Konstanten \\
\href{https://arxiv.org/abs/2604.07130}{2604.07130} & Okuyama \& Ohzeki, 1D LR-Spinglas auf der Nishimori-Linie & 2026 & Offenes Fenster 3/2 ≤ α ≤ 2 \\
\href{https://arxiv.org/abs/2403.00421}{2403.00421} & „Monte Carlo based techniques for quantum magnets with long-range interactions'' & 2024 & Die Pipeline, die GZL im Quantenfall ersetzt \\
\end{longtable}

\textbf{Bewertung.}

\begin{longtable}[]{@{}
  >{\raggedright\arraybackslash}p{(\columnwidth - 4\tabcolsep) * \real{0.1851}}
  >{\raggedright\arraybackslash}p{(\columnwidth - 4\tabcolsep) * \real{0.1140}}
  >{\raggedright\arraybackslash}p{(\columnwidth - 4\tabcolsep) * \real{0.6709}}@{}}
\toprule\noalign{}
\begin{minipage}[b]{\linewidth}\raggedright
Kriterium
\end{minipage} & \begin{minipage}[b]{\linewidth}\raggedright
Punkte
\end{minipage} & \begin{minipage}[b]{\linewidth}\raggedright
Begründung
\end{minipage} \\
\midrule\noalign{}
\endhead
\bottomrule\noalign{}
\endlastfoot
Neuheit & 4 & Keine aktuellen klassischen LR-Reihen gefunden. Alte Arbeiten sind ungeprüft. \\
Passung & 3 & \texttt{evaluate\_graph} passt direkt. Korpus, Vertexgewichte und Hardcore-Abbildung sind Eigenbau. \\
Hebel & 4 & Deterministische, nicht-universelle Größen über ganze σ-Bereiche, dazu die ganze Brillouin-Zone \\
Machbarkeit & 2 & Graphenenumeration und Reihenanalyse sind Monate Arbeit. \\
Wow & 3 & β\textsubscript{c}(σ)-Kurven gegen MC-Punkte, aber abstrakt \\
Anschluss & 5 & Mitten in einem aktiven Streit zwischen drei Communities \\
\end{longtable}

\subsection{Idee 3: Modulare Graphfunktionen als duale Graph-Zetas, GZL als Numerik-Labor der Stringtheorie}

\textbf{Pitch.} Modulare Graphfunktionen (MGFs) aus Genus-1-Stringamplituden sind für planare Graphen exakt GZL-Graph-Zetas des dualen Graphen auf dem Impulsgitter Λ\_τ = ℤ + τℤ. Die Physik der langreichweitigen Quantenmagnete und die Welt der Stringamplituden teilen damit dasselbe numerische Objekt.

\textbf{Das Problem im Zielfeld.} MGFs sind „non-holomorphic modular functions associated with Feynman graphs for a conformal scalar field theory on a two-dimensional torus'' \textbf{{[}belegt{]}} (\href{https://arxiv.org/abs/1512.06779}{arXiv:1512.06779}). Der Stand der Werkzeuge:

\begin{itemize}
\tightlist
\item
  Alle Identitäten bis Gewicht 6 und alle dihedralen Identitäten bei Gewicht 7 sind bewiesen (\href{https://arxiv.org/abs/1608.04393}{arXiv:1608.04393}).
\item
  Ein Mathematica-Paket überführt MGFs in iterierte Eisenstein-Integrale, deckt aber nur Topologien \textbf{bis vier Vertices} ab (\href{https://arxiv.org/abs/2502.05531}{arXiv:2502.05531}).
\item
  Elliptische MGFs werden über äquivariante iterierte Integrale gelöst (\href{https://arxiv.org/abs/2511.15883}{arXiv:2511.15883}).
\item
  Objekte mit komplexen Indizes tauchen neu auf (\href{https://arxiv.org/abs/2512.21413}{arXiv:2512.21413}).
\end{itemize}

Numerische Werte bei beliebigem innerem τ für größere planare Graphen und für nicht-ganzzahlige Exponenten gibt es als Werkzeug nicht allgemein.

\textbf{Abbildung auf GZL {[}abgeleitet, numerisch gestützt{]}.} Ein MGF ist C\_Γ(τ) = Σ\_\{p\_e ∈ Λ\_τ∖0\} Π\textsubscript{v} δ(Σ ±p\_e) Π\textsubscript{e} (τ₂/(π\textbar p\_e\textbar²))\textsuperscript{a\_e}. Das Impulserhaltungsgitter ist das Zyklengitter von Γ. Nach Whitney lässt es sich genau dann als Differenzen von Vertexvariablen schreiben, wenn Γ planar ist, und diese Vertices sind die Flächen von Γ, also die Knoten des Duals Γ*. Daraus folgt:

C\_Γ(τ) = (τ₂/π)\textsuperscript{Σa\_e} · ζ\_\{Λ\_τ, Γ*\}(k = 0) mit ν\textsubscript{e} = 2a\_e.

Die Zuordnung im Einzelnen:

\begin{itemize}
\tightlist
\item
  Dihedrale („Banana''-)MGFs werden zu Kreis-Zetas (\texttt{zeta\_circle}).
\item
  Trihedrale werden zu Theta-artigen SP-Graphen.
\item
  Der selbstduale Tetraeder K₄ läuft über das Tensornetz.
\item
  Die Kante mit Impuls, Z\textsubscript{Λ\_τ,2a}(k), ist der Baustein aus Green-Funktion bzw. Kronecker-Eisenstein-Reihe.
\item
  GZLs externes k entspräche dem Torus-Punkt z elliptischer MGFs \textbf{{[}spekulativ, nur im dihedralen Fall plausibilisiert{]}}.
\end{itemize}

\begin{verbatim}
A = np.array([[1, tau.real], [0, tau.imag]])                 # |A(m,n)| = |m + n*tau|
C_abc = (tau.imag/np.pi)**(a+b+c) * gzl.zeta_circle([2*a, 2*b, 2*c], A)   # dihedral, 2 Vertices
C_K4  = (tau.imag/np.pi)**12 * gzl.evaluate_graph(K4, np.full(6, 4.0), A, n_points=64)  # tetraedrisch
\end{verbatim}

\textbf{Pilotergebnisse {[}Pilot{]}} (\texttt{mgf\_test.py}, \texttt{mgf\_test2.py}). Die Nachprüfung ohne GZL bestätigt beide Identitäten auf 10⁻¹⁴ bis 10⁻¹⁵ (Abschnitt 7.4):

\begin{itemize}
\tightlist
\item
  \textbf{C₂,₂,₁₊ε/₂ → (2/5)E₅ + ζ(5)/30.} Nach Richardson-Extrapolation in ε liegt der relative Fehler bei \textbf{1,1--1,2×10⁻⁸} an allen vier getesteten τ (i, e\textsuperscript{iπ/3}, 0,21 + 1,37i, −0,4 + 0,95i).
\item
  Die \textbf{Zagier-Identität C₁,₁,₁ = E₃ + ζ(3)} ist auf etwa 9×10⁻⁷ reproduziert.
\item
  Jede Kreis-Auswertung braucht \textasciitilde0,1 s.
\item
  Die tetraedrische MGF (alle ν = 4) bei τ = i konvergiert mit n\textsubscript{points} = 32/48/64 gegen 2,29641310909 / \ldots926 / \ldots927, in 0,1--0,3 s.
\item
  Die Zuschreibung von C₂,₂,₁ an D'Hoker--Green--Vanhove ist \emph{nicht verifiziert}.
\end{itemize}

\textbf{Was neu herauskäme.}

\begin{itemize}
\tightlist
\item
  Double-Precision-Werte planarer MGFs mit \textbf{≥ 5 Vertices} an beliebigem innerem τ, jenseits des 4-Vertex-Pakets.
\item
  \textbf{Verallgemeinerte MGFs mit reellen Exponenten} C\textsubscript{s₁,s₂,s₃}(τ).
\item
  Elliptische MGFs auf einem vollen z-Gitter per FFT (spekulativ).
\item
  Ein neuer Satz über modulare Funktionen: Mit dem Wörterbuch aus Idee 1 wandert \textbf{das Minimum von C\textsubscript{s,s,s}(τ) über der Fundamentaldomäne bei s* = 1,95918 vom hexagonalen zum quadratischen Punkt und für s \textgreater{} 4,303 zu einem Rechteckpunkt} {[}nachgeprüft, Abschnitt 7{]}. Nach den Notizen hat diese Frage noch niemand gestellt.
\end{itemize}

\textbf{Erster Meilenstein (3--5 Wochen).} Die dihedralen, trihedralen und tetraedrischen Identitäten an vielen τ systematisch reproduzieren. Die inhomogenen Laplace-Gleichungen numerisch per finiter Differenzen in τ prüfen.

\textbf{Ausbaupfad.}

\begin{itemize}
\tightlist
\item
  \textbf{Monat 2--4:} Tabellen von 5- und 6-Vertex-MGFs. Der Vergleich mit \href{https://arxiv.org/abs/2502.05531}{arXiv:2502.05531} ist nur bis 4 Vertices möglich.
\item
  \textbf{Monat 4--6:} C\textsubscript{s,s,s} für reelle s, Minimierungsstudie und Anschluss an Idee 1.
\item
  \textbf{Monat 6--12:} elliptische MGFs gegen publizierte Werte prüfen. Optional eine Multiprecision-Portierung (mpmath oder Arb) der Kreis-Zeta für PSLQ.
\end{itemize}

\textbf{Vorwissen und Einstieg.} Modulformen, Eisenstein-Reihen, Grundlagen der Genus-1-Stringamplituden, Planarität und Dualität von Graphen. Einstieg: \href{https://arxiv.org/abs/1512.06779}{arXiv:1512.06779}, \href{https://arxiv.org/abs/1608.04393}{arXiv:1608.04393}, \href{https://arxiv.org/abs/2502.05531}{arXiv:2502.05531}.

\textbf{Risiken.}

\begin{enumerate}
\def\labelenumi{\arabic{enumi}.}
\tightlist
\item
  Die MGF-Community hat exakte Methoden. GZL erschließt dort \textbf{kein blockiertes Problem}, sondern liefert Kreuzchecks und Erweiterungen.
\item
  Die meisten stringrelevanten MGFs haben Kanten mit a\textsubscript{e} = 1, also ν = d = 2. Das liegt außerhalb der GZL-Domäne und ist nur per ε-Extrapolation mit \textasciitilde10⁻⁸ erreichbar.
\item
  Nicht-planare MGFs (K₃,₃-Topologien) sind keine GZL-Objekte.
\item
  Modulare Graph\emph{formen} mit a ≠ b brauchen anisotrope Kerne.
\item
  Mit doppelter Genauigkeit ist PSLQ nur für winzige Basen möglich.
\end{enumerate}

\textbf{Scoop-Risiko: gering.} Kein Bezug zur Roadmap der Autoren.

\textbf{Schlüsselpapers.}

\begin{longtable}[]{@{}
  >{\raggedright\arraybackslash}p{(\columnwidth - 6\tabcolsep) * \real{0.1032}}
  >{\raggedright\arraybackslash}p{(\columnwidth - 6\tabcolsep) * \real{0.4037}}
  >{\raggedright\arraybackslash}p{(\columnwidth - 6\tabcolsep) * \real{0.0593}}
  >{\raggedright\arraybackslash}p{(\columnwidth - 6\tabcolsep) * \real{0.4037}}@{}}
\toprule\noalign{}
\begin{minipage}[b]{\linewidth}\raggedright
arXiv
\end{minipage} & \begin{minipage}[b]{\linewidth}\raggedright
Titel / Inhalt
\end{minipage} & \begin{minipage}[b]{\linewidth}\raggedright
Jahr
\end{minipage} & \begin{minipage}[b]{\linewidth}\raggedright
Relevanz
\end{minipage} \\
\midrule\noalign{}
\endhead
\bottomrule\noalign{}
\endlastfoot
\href{https://arxiv.org/abs/1512.06779}{1512.06779} & Modulare Graphfunktionen und single-valued elliptische Polylogarithmen & 2015 & Definition. Impuls an nur zwei Vertices entspricht GZLs Zweiterminal-Struktur \\
\href{https://arxiv.org/abs/1608.04393}{1608.04393} & D'Hoker \& Kaidi, alle Identitäten bis Gewicht 6 & 2016 & Prüfstein für die Numerik \\
\href{https://arxiv.org/abs/2502.05531}{2502.05531} & Claasen \& Doroudiani, iterierte Eisenstein-Integrale, Paket bis 4 Vertices & 2025 & Stand der Technik, Grenze bei 4 Vertices \\
\href{https://arxiv.org/abs/2511.15883}{2511.15883} & Schlotterer, Sohnle, Tao, elliptische MGFs über äquivariante iterierte Integrale & 2025 & Ziel für k-aufgelöste GZL-Werte \\
\href{https://arxiv.org/abs/2603.20550}{2603.20550} & Dobrowolski, symmetrisierte Mordell--Tornheim-Zeta & 2026 & 1D-Kreis-Zeta = ζ̄\_n, Validierungsbank \\
\href{https://arxiv.org/abs/2512.21413}{2512.21413} & Fedosova \& Klinger-Logan, Divisor-Summen mit komplexen Indizes & 2025 & Motivation für nicht-ganzzahlige Exponenten \\
\end{longtable}

\textbf{Bewertung.}

\begin{longtable}[]{@{}
  >{\raggedright\arraybackslash}p{(\columnwidth - 4\tabcolsep) * \real{0.1851}}
  >{\raggedright\arraybackslash}p{(\columnwidth - 4\tabcolsep) * \real{0.1140}}
  >{\raggedright\arraybackslash}p{(\columnwidth - 4\tabcolsep) * \real{0.6709}}@{}}
\toprule\noalign{}
\begin{minipage}[b]{\linewidth}\raggedright
Kriterium
\end{minipage} & \begin{minipage}[b]{\linewidth}\raggedright
Punkte
\end{minipage} & \begin{minipage}[b]{\linewidth}\raggedright
Begründung
\end{minipage} \\
\midrule\noalign{}
\endhead
\bottomrule\noalign{}
\endlastfoot
Neuheit & 5 & Das Wörterbuch MGF ↔ duale Graph-Zeta ist in den Quellen nirgends ausgesprochen. \\
Passung & 3 & Planar und ν \textgreater{} d funktioniert. Die typischen a = 1-Kanten sind Grenzfälle. \\
Hebel & 2 & Exakte Konkurrenzmethoden existieren. \\
Machbarkeit & 4 & Die Piloten laufen, das Wörterbuch ist einfach. \\
Wow & 4 & „Magnetismus-Software berechnet Stringamplituden'' ist eine starke Geschichte. \\
Anschluss & 3 & Interesse als Kreuzcheck. Neue s*-Aussage für math.NT \\
\end{longtable}

\subsection{Idee 4: Simulator-Material-Zwillinge mit vollem Schwanz -- Rydberg-vdW-Arrays und Tm-Dreiecksmagnete}

\textbf{Pitch.} Eine Referenz im thermodynamischen Limes für Rydberg-Quantenprozessoren mit 256 Qubits: volle 1qp-Dispersion, Gap und kritischer Punkt mit dem \emph{kompletten} 1/r⁶-Schwanz. Dazu wird exakt beziffert, wie stark der Schwanz die Abbildung „Simulator ↔ Material'' für TmMgGaO₄ und NaTmSe₂ verschiebt.

\textbf{Das Problem im Zielfeld {[}belegt{]}.}

\begin{itemize}
\tightlist
\item
  Pasqal fährt das 2D-TFIM auf 16×16 Atomen mit U(R/r)⁶. Den kritischen Punkt h*/J ≈ 2,6 für L = 16 übernimmt die Arbeit aus QMC. Tensornetze verlieren bei tJ ≳ 1 die Kontrolle, und Spektren werden nicht gezeigt (\href{https://arxiv.org/abs/2608.07178}{arXiv:2608.07178}).
\item
  Die „one-to-one''-Simulation von TmMgGaO₄ auf einem Dreiecks-Rhombus mit 256 Qubits bildet ein J₁--J₂-Modell mit J₂ ≈ 0,05 J₁ ab. Sie behauptet nur „almost the same neighbour dependence''. Die MPS-Vergleichsrechnungen dauerten etwa zwei Wochen gegenüber etwa einem Tag auf dem QPU (\href{https://arxiv.org/abs/2603.20372}{arXiv:2603.20372}).
\item
  NaTmSe₂ realisiert das TFIM mit quantitativ bestimmten Austauschparametern und Neutronendaten (\href{https://arxiv.org/abs/2505.09884}{arXiv:2505.09884}).
\item
  Für KTmSe₂ zeigte die GZL-Physik-Arbeit, dass „J₁ + voller Dipol'' die INS-Dispersion am besten beschreibt (\href{https://arxiv.org/html/2609.18761}{arXiv:2609.18761}, Daten aus \href{https://arxiv.org/abs/2306.03544}{arXiv:2306.03544}).
\item
  Keines dieser Experimente wird mit einer konvergierten Rechnung im thermodynamischen Limes über die volle Brillouin-Zone mit komplettem Schwanz verglichen.
\end{itemize}

\textbf{Abbildung auf GZL {[}abgeleitet{]}.}

\begin{itemize}
\tightlist
\item
  \textbf{Gitter:} \texttt{triangular} (auch \texttt{square} und \texttt{chain}).
\item
  \textbf{Korpus:} \texttt{tfim1qp} bzw. \texttt{tfim0qp}, AFM-Vorzeichen.
\item
  \textbf{Kern:} Interaction = a(x) + C₆\textbar x\textbar{}\textsuperscript{−6} für Rydberg. Für Tm-Materialien mit Momenten ∥ c ist die Dipol-zz-Kopplung in der Ebene exakt isotrop, also a(x) = J₁, J₂ kompakt plus D\textbar x\textbar{}\textsuperscript{−3}.
\item
  \textbf{Observablen:} 1qp-Dispersion ω(k) über die ganze Brillouin-Zone, Gap am K-Punkt, λ\textsubscript{c} und zν per DlogPadé.
\item
  Der vdW-Schwanz liefert auf dem Dreiecksgitter J₂/J₁ = 1/27 ≈ 0,037 und J₃/J₁ = 1/64. Das Material hat J₂ ≈ 0,05 J₁ und kein J₃. Die Differenz lässt sich genau als kompakte Korrektur a(x) einspeisen.
\end{itemize}

\begin{verbatim}
# Rydberg-Referenz (Z2-symmetrische Detuning-Linie), Pseudokonfiguration
gzl series --corpus tfim1qp --A triangular --nu 6 --n-points 48 --order-max 11
# Materialvariante: Interaction.from_shells(...)  (J1, Delta J2 kompakt) + b*|x|^-3 bzw. |x|^-6
\end{verbatim}

\textbf{Was neu herauskäme.}

\begin{itemize}
\tightlist
\item
  Die ersten Spektren im thermodynamischen Limes mit vollem vdW-Schwanz für die Gitter der 256-Qubit-Maschinen.
\item
  Eine quantifizierte „Schwanzkorrektur'' der TmMgGaO₄-Abbildung: Um wie viel verschieben 1/27 statt 0,05 und der J₃-Schwanz Gap und kritisches Feld?
\item
  Eine Modellwahl J₁--J₂ gegen J₁ + Dipol für NaTmSe₂ nach dem KTmSe₂-Muster, samt Abstand zum Quantenkritischen Punkt.
\end{itemize}

\textbf{Erster Meilenstein (3--6 Wochen).}

\begin{enumerate}
\def\labelenumi{\arabic{enumi}.}
\tightlist
\item
  Das 2D-KTmSe₂-Ergebnis aus \href{https://arxiv.org/html/2609.18761}{arXiv:2609.18761} reproduzieren (Validierung).
\item
  Reines vdW-Modell gegen J₁--J₂-Modell auf dem Dreiecksgitter vergleichen: Dispersion, Gap und λ\textsubscript{c}.
\end{enumerate}

\textbf{Ausbaupfad.}

\begin{itemize}
\tightlist
\item
  \textbf{Monat 2--4:} Vergleich mit Pasqal (h*/J ≈ 2,6) nach Klärung der Vorzeichen- und Untergitter-Konvention. Die Notizen markieren diesen Punkt als \emph{nicht verifiziert}.
\item
  \textbf{Monat 4--6:} NaTmSe₂ und 1D-Rydberg-Ketten für die Caltech-Messungen des dynamischen Strukturfaktors (\href{https://arxiv.org/abs/2601.16275}{arXiv:2601.16275}).
\item
  \textbf{Später:} 3D-Stapelung und Spektralgewichte, sobald die Roadmap-Features vorliegen.
\end{itemize}

\textbf{Vorwissen und Einstieg.} pCUT und Linked-Cluster-Grundlagen, DlogPadé, Rydberg-Hamiltonian, Kristallfeld-Physik nicht-Kramers-Ionen. Einstieg: \href{https://arxiv.org/html/2609.18761}{arXiv:2609.18761}, \href{https://arxiv.org/abs/2603.20372}{arXiv:2603.20372}.

\textbf{Risiken.}

\begin{enumerate}
\def\labelenumi{\arabic{enumi}.}
\tightlist
\item
  Der mitgelieferte Korpus gilt nur auf der \textbf{Z₂-symmetrischen Linie} δ = Σ\textsubscript{j} U\textsubscript{ij}/2. Ein Longitudinalfeld, also das volle (Ω, δ)-Diagramm, braucht einen neuen Korpus.
\item
  pCUT aus dem Transversalfeld-Limes gilt nur in der paramagnetischen Phase. TmMgGaO₄ ordnet bei tiefer Temperatur.
\item
  Ordnung 11 begrenzt die Präzision nahe λ\textsubscript{c}.
\item
  Die Stapelung in 3D braucht anisotrope Dipolkerne.
\end{enumerate}

\textbf{Scoop-Risiko: hoch.} Die Autoren nennen Rydberg-vdW-LRTFIM und DSF-Vergleiche ausdrücklich als nächstes Ziel und haben KTmSe₂ selbst bearbeitet. Empfehlung: Kooperation suchen oder sich auf das konkrete Paar TmMgGaO₄-Simulator und NaTmSe₂ konzentrieren.

\textbf{Schlüsselpapers.}

\begin{longtable}[]{@{}
  >{\raggedright\arraybackslash}p{(\columnwidth - 6\tabcolsep) * \real{0.1147}}
  >{\raggedright\arraybackslash}p{(\columnwidth - 6\tabcolsep) * \real{0.4486}}
  >{\raggedright\arraybackslash}p{(\columnwidth - 6\tabcolsep) * \real{0.0659}}
  >{\raggedright\arraybackslash}p{(\columnwidth - 6\tabcolsep) * \real{0.3408}}@{}}
\toprule\noalign{}
\begin{minipage}[b]{\linewidth}\raggedright
arXiv
\end{minipage} & \begin{minipage}[b]{\linewidth}\raggedright
Titel / Inhalt
\end{minipage} & \begin{minipage}[b]{\linewidth}\raggedright
Jahr
\end{minipage} & \begin{minipage}[b]{\linewidth}\raggedright
Relevanz
\end{minipage} \\
\midrule\noalign{}
\endhead
\bottomrule\noalign{}
\endlastfoot
\href{https://arxiv.org/abs/2603.20372}{2603.20372} & 256-Qubit-„one-to-one''-Simulation von TmMgGaO₄ & 2026 & Zielexperiment für die Schwanzkorrektur \\
\href{https://arxiv.org/abs/2608.07178}{2608.07178} & Pasqal/Browaeys, 2D-TFIM auf 16×16 mit vdW & 2026 & Kritischer Punkt aus QMC bei L = 16 \\
\href{https://arxiv.org/abs/2505.09884}{2505.09884} & NaTmSe₂ als TFIM-Realisierung & 2025 & Nächstes Material nach KTmSe₂ \\
\href{https://arxiv.org/abs/2306.03544}{2306.03544} & KTmSe₂-INS, dispersiver Kristallfeldzweig & 2023 & Validierungsdaten \\
\href{https://arxiv.org/abs/2609.18761}{2609.18761} & Duft et al., „Exact and fast series expansions for quantum models with long-range interactions'' & 2026 & Methode und KTmSe₂-Präzedenz \\
\href{https://arxiv.org/abs/2601.16275}{2601.16275} & Caltech, DSF einer Rydberg-Kette am kritischen Punkt & 2026 & 1D-Vergleichsziel \\
\end{longtable}

\textbf{Bewertung.}

\begin{longtable}[]{@{}
  >{\raggedright\arraybackslash}p{(\columnwidth - 4\tabcolsep) * \real{0.1851}}
  >{\raggedright\arraybackslash}p{(\columnwidth - 4\tabcolsep) * \real{0.1140}}
  >{\raggedright\arraybackslash}p{(\columnwidth - 4\tabcolsep) * \real{0.6709}}@{}}
\toprule\noalign{}
\begin{minipage}[b]{\linewidth}\raggedright
Kriterium
\end{minipage} & \begin{minipage}[b]{\linewidth}\raggedright
Punkte
\end{minipage} & \begin{minipage}[b]{\linewidth}\raggedright
Begründung
\end{minipage} \\
\midrule\noalign{}
\endhead
\bottomrule\noalign{}
\endlastfoot
Neuheit & 3 & Neue Systeme, aber dieselbe Methode wie bei den Autoren \\
Passung & 4 & Läuft heute, auf der Z₂-Linie und im 2D-Dipolmodell \\
Hebel & 4 & Voller Schwanz und volle Brillouin-Zone im thermodynamischen Limes statt DMRG/QMC bei L = 16 \\
Machbarkeit & 4 & Mitgelieferte Korpora, das Kern-API reicht aus. \\
Wow & 5 & Direkter Vergleich mit Experiment und Material \\
Anschluss & 4 & Relevant für Rydberg- und Neutronen-Community. Kooperation ist naheliegend. \\
\end{longtable}

\subsection{Idee 5: Quanten-Sak-Test und Log-Korrekturen im LRTFIM für d = 1, 2, 3}

\textbf{Pitch.} Mit rauschfreien Reihen auf einem sehr dichten σ-Gitter prüfen, ob die Exponenten des LRTFIM ihren kurzreichweitigen Wert bei σ = 2 − η\textsubscript{SR} (Sak) oder erst bei σ = 2 erreichen, wie es das RG-Papier von 2026 behauptet. Dazu die multiplikativen Logarithmen an der oberen kritischen Grenze σ\textsubscript{uc} = 2d/3 und im 3D-Kubikgitter vermessen.

\textbf{Das Problem im Zielfeld {[}belegt{]}.} \href{https://arxiv.org/abs/2606.22407}{arXiv:2606.22407} schlägt ein (d, σ)-Universalitätsdiagramm für Quanten-O(n) vor, mit σ* = 2 und Vorhersagen für ν, η\_ω und η\textsubscript{k}. Klassisch behauptet die Deng-Gruppe dasselbe per MC (\href{https://arxiv.org/abs/2512.04805}{arXiv:2512.04805}). Für die antiferromagnetische LR-Ising-Kette weichen Korrelationsexponent und zentrale Ladung laut NQS bei kleinem Zerfallsexponenten ab, und der Zusammenbruch der konformen Invarianz ist „yet to be determined'' (\href{https://arxiv.org/abs/2308.09709}{arXiv:2308.09709}). Frühere pCUT+MC-Reihen waren verrauscht und nur dünn in σ abgetastet (\href{https://arxiv.org/abs/1802.06684}{arXiv:1802.06684}, \href{https://arxiv.org/abs/2203.08081}{arXiv:2203.08081}).

\textbf{Abbildung auf GZL.} Direkt mit den mitgelieferten Korpora \textbf{{[}belegt{]}}: ν = d + σ auf dem Gitter \texttt{chain}, \texttt{square}, \texttt{triangular} oder \texttt{cubic}, FM- und AFM-Vorzeichen, 1qp-Gap als Observable. Die Reihenanalyse umfasst DlogPadé und Ratio-Methoden mit Log-Termen.

\begin{verbatim}
for sigma in $(seq 1.50 0.005 2.20); do
  gzl series --corpus tfim1qp --A chain --nu $(echo "1+$sigma" | bc) --order-max 11 ...
done   # -> z*nu(sigma); d=1: Sak 7/4 vs 2; d=2: ~1,96 vs 2
\end{verbatim}

\textbf{Was neu herauskäme.}

\begin{itemize}
\tightlist
\item
  Deterministische Kurven zν(σ) in d = 1, 2, 3, direkt gegen die Vorhersagen aus \href{https://arxiv.org/abs/2606.22407}{arXiv:2606.22407} gestellt.
\item
  Die erste quantitative Extraktion multiplikativer Log-Exponenten aus exakten Koeffizienten in d = 2 und 3. In d = 3 gehört die kurzreichweitige Klasse zu 4D-Ising und trägt Logarithmen im gesamten SR-Regime \textbf{{[}abgeleitet{]}}.
\end{itemize}

\textbf{Erster Meilenstein (2--4 Wochen).} Zuerst den Crossover-Abschnitt von \href{https://arxiv.org/html/2609.18761}{arXiv:2609.18761} lesen. Die Notizen konnten dessen Zahlen nicht extrahieren, und davon hängt die Neuheit ab. Danach zν(σ) für die FM-Kette auf dem Gitter σ ∈ {[}1,5; 2,2{]} mit Schrittweite 0,005 berechnen.

\textbf{Ausbaupfad.} Zuerst d = 2 (Quadrat, Dreieck AFM), dann das 3D-Kubikgitter, dann Log-bewusste Extrapolation. Langfristig O(2) und O(3) über XXZ-Korpora, die heute fehlen.

\textbf{Vorwissen und Einstieg.} Kritische Phänomene langreichweitiger Modelle, Reihenanalyse. Einstieg: \href{https://arxiv.org/abs/2203.08081}{arXiv:2203.08081}, \href{https://arxiv.org/abs/2606.22407}{arXiv:2606.22407}.

\textbf{Risiken.}

\begin{enumerate}
\def\labelenumi{\arabic{enumi}.}
\tightlist
\item
  In d = 2 liegen die beiden Kandidaten nur \textbf{\textasciitilde0,04} auseinander. Bei Ordnung 11 und langsamen Crossovers begrenzt der Extrapolationsfehler, nicht das Rauschen.
\item
  Die Autoren haben σ bereits dicht gescannt. Die Neuheit muss aus der expliziten Konfrontation mit RG und MC und aus der Log-Analyse kommen.
\end{enumerate}

\textbf{Scoop-Risiko: mittel--hoch.} Die GZL-Arbeit charakterisiert „the crossover between universality regimes'' bereits.

\textbf{Schlüsselpapers.}

\begin{longtable}[]{@{}
  >{\raggedright\arraybackslash}p{(\columnwidth - 6\tabcolsep) * \real{0.1139}}
  >{\raggedright\arraybackslash}p{(\columnwidth - 6\tabcolsep) * \real{0.4454}}
  >{\raggedright\arraybackslash}p{(\columnwidth - 6\tabcolsep) * \real{0.0654}}
  >{\raggedright\arraybackslash}p{(\columnwidth - 6\tabcolsep) * \real{0.3453}}@{}}
\toprule\noalign{}
\begin{minipage}[b]{\linewidth}\raggedright
arXiv
\end{minipage} & \begin{minipage}[b]{\linewidth}\raggedright
Titel / Inhalt
\end{minipage} & \begin{minipage}[b]{\linewidth}\raggedright
Jahr
\end{minipage} & \begin{minipage}[b]{\linewidth}\raggedright
Relevanz
\end{minipage} \\
\midrule\noalign{}
\endhead
\bottomrule\noalign{}
\endlastfoot
\href{https://arxiv.org/abs/2606.22407}{2606.22407} & „Perturbative Renormalization and Universality Diagram for Long-Range Quantum Criticality'' & 2026 & Zu prüfende Vorhersage σ* = 2 \\
\href{https://arxiv.org/abs/2203.08081}{2203.08081} & „Scaling at quantum phase transitions above the upper critical dimension'' & 2022 & Anker für Exponenten und Log-Korrekturen \\
\href{https://arxiv.org/abs/1802.06684}{1802.06684} & „Quantum criticality of two-dimensional quantum magnets with long-range interactions'' & 2019 & MC-basierte Vorgängerreihen \\
\href{https://arxiv.org/abs/2512.04805}{2512.04805} & Klassisch σ* = 2 per MC & 2025 & Klassisches Gegenstück \\
\href{https://arxiv.org/abs/2308.09709}{2308.09709} & NQS für die AF-LR-Ising-Kette & 2023 & Dokumentierte Uneinigkeit \\
\href{https://arxiv.org/abs/2609.18761}{2609.18761} & GZL-Physik-Arbeit & 2026 & Ausgangspunkt, Überlappung prüfen \\
\end{longtable}

\textbf{Bewertung.}

\begin{longtable}[]{@{}
  >{\raggedright\arraybackslash}p{(\columnwidth - 4\tabcolsep) * \real{0.1851}}
  >{\raggedright\arraybackslash}p{(\columnwidth - 4\tabcolsep) * \real{0.1140}}
  >{\raggedright\arraybackslash}p{(\columnwidth - 4\tabcolsep) * \real{0.6709}}@{}}
\toprule\noalign{}
\begin{minipage}[b]{\linewidth}\raggedright
Kriterium
\end{minipage} & \begin{minipage}[b]{\linewidth}\raggedright
Punkte
\end{minipage} & \begin{minipage}[b]{\linewidth}\raggedright
Begründung
\end{minipage} \\
\midrule\noalign{}
\endhead
\bottomrule\noalign{}
\endlastfoot
Neuheit & 3 & Die Überlappung mit dem σ-Scan der Autoren ist wahrscheinlich. \\
Passung & 5 & Nur mitgelieferte Korpora \\
Hebel & 3 & Das Rauschen verschwindet, das Ordnungslimit bleibt. \\
Machbarkeit & 4 & Rechenzeit in Minuten, die Analyse ist die Arbeit. \\
Wow & 3 & Exponentenkurven \\
Anschluss & 4 & Aktiver Streit \\
\end{longtable}

\begin{center}\rule{0.5\linewidth}{0.5pt}\end{center}

\section{Weitere Ideen: vierzehn Kandidaten mit ehrlichem Erweiterungsbedarf}

N = Neuheit, P = Passung, H = Hebel, M = Machbarkeit, W = Wow, A = Anschlussfähigkeit (je 1--5).

\begin{longtable}[]{@{}
  >{\raggedright\arraybackslash}p{(\columnwidth - 20\tabcolsep) * \real{0.2311}}
  >{\raggedright\arraybackslash}p{(\columnwidth - 20\tabcolsep) * \real{0.0339}}
  >{\raggedright\arraybackslash}p{(\columnwidth - 20\tabcolsep) * \real{0.2311}}
  >{\raggedright\arraybackslash}p{(\columnwidth - 20\tabcolsep) * \real{0.2311}}
  >{\raggedright\arraybackslash}p{(\columnwidth - 20\tabcolsep) * \real{0.0339}}
  >{\raggedright\arraybackslash}p{(\columnwidth - 20\tabcolsep) * \real{0.0339}}
  >{\raggedright\arraybackslash}p{(\columnwidth - 20\tabcolsep) * \real{0.0339}}
  >{\raggedright\arraybackslash}p{(\columnwidth - 20\tabcolsep) * \real{0.0339}}
  >{\raggedright\arraybackslash}p{(\columnwidth - 20\tabcolsep) * \real{0.0339}}
  >{\raggedright\arraybackslash}p{(\columnwidth - 20\tabcolsep) * \real{0.0339}}
  >{\raggedright\arraybackslash}p{(\columnwidth - 20\tabcolsep) * \real{0.0393}}@{}}
\toprule\noalign{}
\begin{minipage}[b]{\linewidth}\raggedright
Idee
\end{minipage} & \begin{minipage}[b]{\linewidth}\raggedright
Spur
\end{minipage} & \begin{minipage}[b]{\linewidth}\raggedright
Kern der Abbildung
\end{minipage} & \begin{minipage}[b]{\linewidth}\raggedright
Was fehlt
\end{minipage} & \begin{minipage}[b]{\linewidth}\raggedright
N
\end{minipage} & \begin{minipage}[b]{\linewidth}\raggedright
P
\end{minipage} & \begin{minipage}[b]{\linewidth}\raggedright
H
\end{minipage} & \begin{minipage}[b]{\linewidth}\raggedright
M
\end{minipage} & \begin{minipage}[b]{\linewidth}\raggedright
W
\end{minipage} & \begin{minipage}[b]{\linewidth}\raggedright
A
\end{minipage} & \begin{minipage}[b]{\linewidth}\raggedright
Scoop
\end{minipage} \\
\midrule\noalign{}
\endhead
\bottomrule\noalign{}
\endlastfoot
ATM-Elastizitätskonstanten, Cauchy-Verletzung C₁₂ − C₄₄, Phonon-Zweitmoment Tr D(q) mit Dreikörperkräften (\href{https://arxiv.org/abs/2504.07338}{2504.07338}, \href{https://arxiv.org/abs/2012.05413}{2012.05413}) & B & Summe von Dreiecks-Zetas, Ableitung nach A, Bloch-Phase auf einer Kante & Fortsetzung für ν ≤ d in Dreiecksblöcken, nur Bravais (kein hcp) & 4 & 2 & 4 & 3 & 3 & 4 & mittel \\
Vierkörper-Dispersion (Dipolring) & B & 4-Zyklus bzw. K₄ & Diagonalen mit positiven Potenzen, anisotrope Kontraktion & 4 & 1 & 3 & 2 & 3 & 3 & mittel \\
1D-Rydberg-Kette am Ising-Punkt für Caltech-CFT-Spektren (\href{https://arxiv.org/abs/2601.16275}{2601.16275}) & A & \texttt{chain}, ν = 6, AFM & nichts auf der Z₂-Linie; tricritisch und E8 brauchen neue Korpora & 3 & 4 & 3 & 4 & 4 & 4 & hoch \\
Langreichweitiger Benchmark für NLCE+QA auf Ionen-QPUs (\href{https://arxiv.org/abs/2605.28599}{2605.28599}) & A & LR-Referenzdispersion für α = 1--3 & nichts & 3 & 5 & 3 & 5 & 2 & 3 & mittel \\
Bilayer- und Dimer-Kerne (r² + d²)\textsuperscript{−α/2} als asymptotische Potenzreihe plus a(x) (\href{https://arxiv.org/abs/2408.13145}{2408.13145}, \href{https://arxiv.org/abs/2605.13969}{2605.13969}) & A & Kerntrick ohne mehratomige Basis & Dimer-Korpus, Konvergenzprüfung & 4 & 3 & 3 & 3 & 2 & 4 & hoch \\
α\textsubscript{c} der LR-XY-Kette (LSW 3,0 gegen ≈ 2,7) und der Spin-1-Kette (QMC 2,49(1)) (\href{https://arxiv.org/abs/2601.20058}{2601.20058}, \href{https://arxiv.org/abs/2604.20831}{2604.20831}, \href{https://arxiv.org/abs/2604.12754}{2604.12754}) & A & Reihen über α & XXZ- und Spin-1-Korpora, gestaffelte Vorzeichen nicht als Kern darstellbar & 2 & 2 & 4 & 2 & 3 & 4 & hoch \\
Dipolares XY/XXZ auf dem Quadratgitter, AFM-Dämpfungsrätsel (\href{https://arxiv.org/abs/2311.11726}{2311.11726}, \href{https://arxiv.org/abs/2605.07685}{2605.07685}) & A & Isotropes 1/r³ bei senkrechter Quantisierungsachse & XXZ-Korpus, gapped Entwicklungslimes & 3 & 2 & 4 & 2 & 4 & 4 & hoch \\
1/S-Spinwellen-Schleifen als Brillouin-Zonen-Faltungen über die Zeta-Algebra & A/B & \texttt{graph\_convolve}-Algebra auf J(k) & Methodenentwicklung {[}spekulativ{]} & 4 & 2 & 3 & 2 & 3 & 3 & gering \\
Magnon-Bindungszustände im 2D-LRTFIM (\href{https://arxiv.org/abs/2512.09037}{2512.09037}) & A & 2qp-Sektor & zwei Impulse, nicht in v1.0.0 & 3 & 1 & 4 & 1 & 4 & 4 & hoch \\
EA-Hochtemperaturreihe des 1D/2D-LR-Spinglases (\href{https://arxiv.org/abs/2604.07130}{2604.07130}) & B & gerade Graphen mit J² ∝ r\textsuperscript{−2σ} & eigener Korpus {[}spekulativ{]} & 4 & 3 & 3 & 2 & 2 & 3 & gering \\
Lévy-Flug-Konstanten: P\textsubscript{n}(0) = Kreis-Zeta/Zⁿ, Torus-Plateaus für FSS oberhalb d\textsubscript{c} (\href{https://arxiv.org/abs/2412.08814}{2412.08814}) & B & Kreis-Zeta bzw. D̂(k) auf dem Gitter & singuläres BZ-Integral, meist genügt EpsteinLib & 2 & 3 & 2 & 4 & 2 & 2 & gering \\
Multipräzisions-Graph-Zeta für PSLQ & B & Port des BZ-Integrals nach mpmath/Arb & Neuimplementierung & 4 & 1 & 3 & 2 & 3 & 3 & gering \\
3D-skalare Gitterstörungstheorie, Vakuumgraphen & B & kompakter Kern plus Potenzschwanz wie bei Lüscher--Weisz & anisotrope Subleading-Terme, IR-Regulator (ν = 1 \textless{} d) & 3 & 1 & 2 & 1 & 2 & 2 & gering \\
Log-Korrekturen im Dicke-Ising-Kontext (\href{https://arxiv.org/abs/2605.27484}{2605.27484}) & A & nur der LR-Ising-Anteil & Photonmode ist all-to-all und keine Gittersumme & 2 & 2 & 2 & 3 & 2 & 3 & hoch \\
\end{longtable}

\begin{center}\rule{0.5\linewidth}{0.5pt}\end{center}

\section{Überraschende Zusammenhänge: eine Zahl, drei Communities}

Der auffälligste Befund dieser Recherche ist eine \textbf{Drei-Welten-Identität} \textbf{{[}abgeleitet, durch Pilot gestützt{]}}. Dieselbe Zahl, die Dreiecks-Graph-Zeta auf einem 2D-Gitter mit Einheitsvolumen, erscheint an drei Stellen:

\begin{itemize}
\tightlist
\item
  als \textbf{Dreikörper-Riesz-Energie} in der Kristallchemie;
\item
  als \textbf{zweischleifige dihedrale MGF} π\textsuperscript{3s}·C\textsubscript{s,s,s}(τ) in der Stringtheorie;
\item
  als \textbf{Einbettungssumme der Zyklen dritter Ordnung} in der Linked-Cluster-Reihe eines langreichweitigen Quantenmagneten.
\end{itemize}

Dahinter steckt eine strukturelle Übereinstimmung: GZLs Darstellung der Kreis-Zeta als Brillouin-Zonen-Integral über Produkte von Epstein-Zetas ist genau die „position-space''-Darstellung von MGFs. Die Torus-Green-Funktion G\textsubscript{a}(z\textbar τ) ist eine Epstein-Zeta am Wellenvektor z (\href{https://arxiv.org/abs/2609.18918}{arXiv:2609.18918}, \href{https://arxiv.org/abs/1512.06779}{arXiv:1512.06779}). In 1D wird dasselbe Objekt zur symmetrisierten Mordell--Tornheim-Zeta aus der Zahlentheorie (\href{https://arxiv.org/abs/2603.20550}{arXiv:2603.20550}). Diese Brücke ist exakt, aber wenig ergiebig, weil 1D bereits über MZVs verstanden ist.

Der \textbf{Sak-Streit wird gleichzeitig in vier Communities ausgetragen}, und keine nutzt Reihen mit exakten Gitterkoeffizienten:

\begin{itemize}
\tightlist
\item
  klassisches MC (\href{https://arxiv.org/abs/2512.04805}{arXiv:2512.04805});
\item
  Perkolation (\href{https://arxiv.org/abs/2608.20750}{arXiv:2608.20750});
\item
  Kontinuums-ε-Entwicklungen (\href{https://arxiv.org/abs/2602.07818}{arXiv:2602.07818}, \href{https://arxiv.org/abs/2608.15120}{arXiv:2608.15120});
\item
  Quanten-RG (\href{https://arxiv.org/abs/2606.22407}{arXiv:2606.22407}).
\end{itemize}

Klassische Hochtemperaturreihe und Quanten-Linked-Cluster-Entwicklung teilen dieselben Gittersummen. GZL ist deshalb der gemeinsame Nenner, der die klassische und die Quantenseite mit \emph{einem} Werkzeug prüfen könnte (Ideen 2 und 5).

Zwei weniger offensichtliche Identitäten gelten exakt \textbf{{[}abgeleitet{]}}. Erstens ist die Rückkehrwahrscheinlichkeit eines Lévy-Flugs auf dem Gitter nach n Schritten P\textsubscript{n}(0) = ζ\textsuperscript{(n)}\emph{Λ(ν,\ldots,ν)/Z}\{Λ,ν\}(0)ⁿ, also eine Kreis-Zeta. Zweitens ist das phononische Zweitmoment Tr D(q) eines Bravais-Kristalls mit ATM-Kräften eine Summe von Dreiecks-Zetas mit Bloch-Phase auf einer Kante. Das ist genau GZLs Modus mit externem Impuls, und damit wird die Gitterdynamik der Kristallchemie zu einem Anwendungsfall der „ganzen Brillouin-Zone per FFT''. Spekulativ ist die Verbindung zu Feynman-Perioden: Das Residuum einer Graph-Zeta am Gesamtpol ν\textbar E\textbar{} = d(\textbar V\textbar−1) sollte die Ortsraum-Periode sein. Das verlangt allerdings ν \textless{} d pro Kante und liegt damit außerhalb von GZL.

Zwei aktuelle Arbeiten liefern schließlich die \textbf{Theorie dafür, warum der volle Schwanz zählt}, und damit ein Verkaufsargument für alle Ideen der Spur A. Das Abschneiden eines r\textsuperscript{−p}-Schwanzes bei Reichweite R verschiebt lokale Observablen um O(R\textsuperscript{−(p−d)}), und diese Schranke ist optimal (\href{https://arxiv.org/abs/2608.15576}{arXiv:2608.15576}). Surrogate aus Exponentialsummen in iMPS verzerren kritische Diagnosen (\href{https://arxiv.org/abs/2606.20522}{arXiv:2606.20522}). Die DMRG-Studie dipolarer XY-Modelle schneidet bei R\textsubscript{max} = W/2 ab (\href{https://arxiv.org/html/2605.07685}{arXiv:2605.07685}).

\begin{center}\rule{0.5\linewidth}{0.5pt}\end{center}

\section{Verworfene Ideen: wo die Brücke nur mit Gewalt hält}

\begin{longtable}[]{@{}
  >{\raggedright\arraybackslash}p{(\columnwidth - 2\tabcolsep) * \real{0.4850}}
  >{\raggedright\arraybackslash}p{(\columnwidth - 2\tabcolsep) * \real{0.4850}}@{}}
\toprule\noalign{}
\begin{minipage}[b]{\linewidth}\raggedright
Idee
\end{minipage} & \begin{minipage}[b]{\linewidth}\raggedright
Warum sie nicht trägt
\end{minipage} \\
\midrule\noalign{}
\endhead
\bottomrule\noalign{}
\endlastfoot
Photonik, Atomarrays, kollektive Lamb-Verschiebung (\href{https://arxiv.org/abs/2512.24596}{2512.24596}, \href{https://arxiv.org/abs/2609.17313}{2609.17313}) & Der Kern ist die komplexe, oszillierende dyadische Green-Funktion e\textsuperscript{ik₀r}/r. GZLs k ist nur eine Bloch-Phase auf einem reellen geraden Kern. Die Bandstruktur ist ohnehin eine Einfachsumme. \\
Lüscher-Zeta und ihre LR-Modifikation (\href{https://arxiv.org/abs/2507.18399}{2507.18399}) & Polkern 1/(p² − q²) mit winkelabhängigen Vertexgewichten, kein Differenzkern \\
QED\_L-Koeffizienten, Madelung-Konstanten (\href{https://arxiv.org/abs/1702.01296}{1702.01296}, \href{https://arxiv.org/abs/2608.10041}{2608.10041}, \href{https://arxiv.org/abs/2503.00977}{2503.00977}) & Einfachsummen. Mehrschleifen-Punktladungssummen faktorisieren. Das richtige Werkzeug ist EpsteinLib, nicht GZL. \\
Ortsraum-Photonpropagator in Gitter-QCD+QED (\href{https://arxiv.org/abs/2603.13086}{2603.13086}) & Die Vertexfunktionen sind MC-Korrelatordaten, keine analytischen Kerne. Außerdem d = 4. \\
Langreichweitige ML-Potenziale (\href{https://arxiv.org/abs/2606.06617}{2606.06617}, \href{https://arxiv.org/abs/2512.18029}{2512.18029}) & Nur Zweikörper-Ewald-Summen. Kein ML-Ansatz braucht nicht-baumartige Graphsummen. \\
Many-Body-Dispersion (MBD, \href{https://arxiv.org/abs/2308.03140}{2308.03140}) & Auf Bravais-Gittern ein 3×3-T(k) pro k: anisotrope \emph{Einfach}summen. Reale Systeme sind mehratomig. \\
LR-Kitaev-Ketten (\href{https://arxiv.org/abs/2412.01076}{2412.01076}, \href{https://arxiv.org/abs/2609.28735}{2609.28735}) & Freie Fermionen, nur eindimensionale Polylog-Transformierte \\
Perkolations-Dreiecksdiagramm nahe s ≈ d & Der Kern wäre das unbekannte kritische τ(x) ≍ ‖x‖\textsuperscript{−d+α}, das kein Eingabekern und nicht in ℓ¹ ist (\href{https://arxiv.org/abs/2404.07276}{2404.07276}). \\
k-Punkt-Baumsummen der Hochdimensions-Perkolation (\href{https://arxiv.org/abs/2607.22387}{2607.22387}) & Die Green-Funktions-Kerne sind nicht summierbar. \\
Feynman-Perioden (\href{https://arxiv.org/abs/2607.25595}{2607.25595}) & ν \textless{} d pro Kante. Graphical functions und HyperInt sind weit überlegen. \\
Epidemien, Ausbreitungskerne & Nur Normierungen niedriger Ordnung, per FFT oder im Kontinuum billig. 2025--26 wurde kein Engpass gefunden. \\
Ionen mit modenbasierten Kopplungen (\href{https://arxiv.org/abs/2606.13499}{2606.13499}) & J\textsubscript{ij} ∝ b\textsubscript{ik} b\textsubscript{jk} ist kein Potenzgesetz. \\
1D-Tornheim- und MZV-Summen & Theoretisch bereits exakt bekannt, GZL taugt nur als Prüfbank. \\
PSLQ-Entdeckungen mit GZL, wie es heute ist & Doppelte Genauigkeit reicht nur für winzige Basen. \\
LiHoF₄, Pyrochlor-Spin-Eis, NaGdS₂ & Nicht-Bravais \textbf{und} anisotroper Dipol \textbf{und} ausdrücklich Ziel der Autoren. Erst nach der Roadmap. \\
Finanzen (Ising-Marktmodelle) & In der Recherche fand sich keine Quelle, die eine Gittersumme mit Potenzgesetz als Engpass zeigt. Nicht weiterverfolgt. \\
\end{longtable}

\begin{center}\rule{0.5\linewidth}{0.5pt}\end{center}

\section{Empfehlung: morgen mit Idee 1 anfangen}

\textbf{Starte mit „Optimale Gitter für Mehrkörper-Potenzgesetzenergien''.} Die Idee hat die beste Kombination aus Neuheit und sofortiger Machbarkeit:

\begin{itemize}
\tightlist
\item
  Sie läuft mit GZL 1.0.0 ohne jede Erweiterung.
\item
  Die Piloten liegen schon vor.
\item
  Sie verbindet drei Welten (Gitteroptimierung, Kristallchemie, modulare Graphfunktionen).
\item
  Das Ergebnis ist ein zeigbares Phasendiagramm.
\item
  Sie berührt die Roadmap der Autoren nur am Rand.
\end{itemize}

Idee 2 ist der natürliche zweite, langfristige Strang mit dem größten Hebel. Die Infrastruktur für Graphenenumeration und Reihenanalyse, die man dafür baut, nützt später auch den Ideen 4 und 5.

\textbf{Schritt 1: Piloten härten (Woche 1--2). Erledigt, siehe Abschnitt 7.} In \texttt{research/pilots/} \texttt{opt\_test.py}, \texttt{opt2.py} und \texttt{opt3.py} erneut laufen lassen (\texttt{mgf\_test2.py} nur aus diesem Verzeichnis). Dann ν* ≈ 3,91836 und die Barriere bei \textasciitilde72° mit variiertem n\textsubscript{points}, mit \texttt{evaluate\_graph} statt \texttt{zeta\_circle} und mit Brute-Force-Summen bei drei ν-Werten absichern. Die Fundamentaldomäne bis Im τ ≈ 5 und mit dichteren Starts absuchen. Ergebnis: eine belastbare 2D-Karte E₃(τ; ν).

\textbf{Schritt 2: globale 3D-Suche (Woche 3--8).} Die 3D-Gitter mit Einheitsvolumen über reduzierte Gram-Matrizen parametrisieren (5 Parameter) und die Dreiecks-Zeta für ν ∈ {[}3,2; 12{]} global minimieren. Ziel ist, festzustellen, ob BCC das globale Minimum ist oder ob ein anderes Gitter gewinnt. Parallel C\textsubscript{s,s,s}(τ) für reelle s gegen die Laplace-Gleichungen der ganzzahligen Fälle prüfen (Brücke zu Idee 3).

\textbf{Schritt 3: Neuheit sichern und Kontakt aufnehmen (parallel, ab Woche 2).} Die offenen Fragen von Bétermin und die Arbeiten von Luo und Wei auf Mehrkörperaussagen durchsehen. Danach Buchheit bzw. den Schwerdtfeger-Kreis kontaktieren: GZL kommt von ihnen, und die offene BCC-Stabilitätsfrage aus \href{https://arxiv.org/abs/2504.07338}{arXiv:2504.07338} ist ihr Thema. So wird das Scoop-Risiko zur Kooperation, und die für ATM nötige Fortsetzung für ν ≤ d kann gemeinsam in GZL einfließen.

\textbf{Schluss.} Durch die Recherche verschiebt sich der Blick: Das eigentliche Potenzial von GZL liegt weniger in „mehr Quantenmagnetismus'' als darin, dass GZL als \textbf{erstes allgemeines Werkzeug für Graph-strukturierte Potenzgesetzsummen} Gebiete verbindet, die dieselben Zahlen unter verschiedenen Namen berechnen. In Spur A haben die Autoren fast jede naheliegende Richtung öffentlich beansprucht. Neuheit entsteht dort nur über konkrete Experimentpaare oder Kooperation. In Spur B ist das Feld leer: Seit dem Release am 25.09.2026 gibt es keine externe Nutzung, und klassische LR-Reihen, Mehrkörper-Gitteroptimierung und das MGF-Wörterbuch tauchen in keiner Roadmap auf. Die kritischste Erweiterung für den Transfer ist dabei nicht die Anisotropie, sondern die \textbf{meromorphe Fortsetzung für Nicht-Brücken-Blöcke mit ν ≤ d}. Sie würde ATM-Kristallphysik und stringrelevante MGFs mit a = 1-Kanten zugleich erschließen.

\begin{center}\rule{0.5\linewidth}{0.5pt}\end{center}

\section{Nachprüfung der Pilotrechnungen: bestätigt, präzisiert und um eine zweite Phase ergänzt}

Die Pilotrechnungen aus den Ideen 1 und 3 wurden nachträglich unabhängig nachgerechnet. Alle Skripte, Rohdaten und Logs liegen in \texttt{research/verification/} und laufen mit \texttt{gzl\ 1.0.0} und \texttt{epsteinlib}. Die Kennzeichnung \textbf{{[}nachgeprüft{]}} heißt: mit zwei unabhängigen Verfahren übereinstimmend berechnet. Das ersetzt keine Begutachtung.

\textbf{Methode.} Als Referenz dient eine eigene Implementierung ohne GZL-Code (\texttt{lattice\_sums.py}). Die Dreieckssumme T\_ν(Λ) = Σ'\textsubscript{x,y} \textbar x\textbar{}\textsuperscript{\{−ν\}\textbar y\textbar{}}\{−ν\}\textbar x−y\textbar{}\textsuperscript{−ν} wird direkt über alle Gitterpunkte in einem Würfel {[}−R, R{]}\textsuperscript{d} summiert. Die innere Faltung läuft per FFT. Danach wird in R extrapoliert (Richardson mit dem bekannten Abschneidefehler \textasciitilde R\textsuperscript{d−2ν}). Die FFT-Version stimmt mit einer naiven Doppelschleife auf Maschinengenauigkeit überein. Für 2D-Gitter ist die Referenz sogar schneller als GZL (0,04 s gegenüber 0,1 s pro Gitter).

\subsection{GZL ist korrekt (V1)}

In 55 Testfällen wurde \texttt{gzl.zeta\_circle} mit der Referenz verglichen: fünf 2D-Gitter und vier 3D-Gitter (FCC, BCC, SC, triklin), mit ν von 2,5 bis 9 und mit ungleichen Exponenten. \textbf{Die größte relative Abweichung beträgt 1,9·10⁻¹⁰, typisch sind 10⁻¹⁵} \textbf{{[}nachgeprüft{]}}. Die größeren Abweichungen treten nur nahe ν = d auf, wo die Extrapolation der Referenz schwieriger wird. \texttt{evaluate\_graph} liefert für das Dreieck dasselbe wie \texttt{zeta\_circle}, bis auf 3·10⁻¹⁶. In 3D braucht GZL 1--160 s pro Gitter, die FFT-Referenz etwa 1 s. Für die globale Suche wurde deshalb die Referenz verwendet und GZL als Kontrolle.

\subsection{2D: drei Phasen statt zwei (V2)}

\begin{longtable}[]{@{}
  >{\raggedright\arraybackslash}p{(\columnwidth - 6\tabcolsep) * \real{0.0801}}
  >{\raggedright\arraybackslash}p{(\columnwidth - 6\tabcolsep) * \real{0.5457}}
  >{\raggedright\arraybackslash}p{(\columnwidth - 6\tabcolsep) * \real{0.1721}}
  >{\raggedright\arraybackslash}p{(\columnwidth - 6\tabcolsep) * \real{0.1721}}@{}}
\toprule\noalign{}
\begin{minipage}[b]{\linewidth}\raggedright
ν
\end{minipage} & \begin{minipage}[b]{\linewidth}\raggedright
globales Minimum (Suche über die ganze Fundamentaldomäne)
\end{minipage} & \begin{minipage}[b]{\linewidth}\raggedright
T hexagonal
\end{minipage} & \begin{minipage}[b]{\linewidth}\raggedright
T Quadrat
\end{minipage} \\
\midrule\noalign{}
\endhead
\bottomrule\noalign{}
\endlastfoot
2,2 & hexagonal & 29,7948755694 & 30,2601286381 \\
3,0 & hexagonal & 13,3917941103 & 13,6528937154 \\
3,9 & hexagonal & 7,9670208300 & 7,9723204026 \\
3,95 & Quadrat & 7,7829092994 & 7,7737904082 \\
6,0 & Quadrat & 3,7457712638 & 3,2499743124 \\
8,0 & Quadrat & 2,2280735200 & 1,5524012126 \\
10 & \textbf{Rechteck, b/a = 1,0998} (T = 0,7597084) & 1,4063994404 & 0,7622381339 \\
20 & \textbf{Rechteck, b/a = 1,1941} (T = 0,0212133) & 0,1603699 & 0,0234490 \\
\end{longtable}

\begin{itemize}
\tightlist
\item
  \textbf{Übergang hexagonal → Quadrat bei ν* = 3,9183649026} \textbf{{[}nachgeprüft{]}}. Referenz und GZL liefern denselben Wert auf zehn Stellen. Das bestätigt den Pilotwert 3,91836.
\item
  \textbf{Der Übergang ist erster Ordnung} \textbf{{[}nachgeprüft{]}}. Die Hesse-Matrix zeigt: Das Quadratgitter ist ab ν = 3,8063 ein lokales Minimum, das Hexagonalgitter nur bis ν = 4,2784. Dazwischen sind beide lokal stabil. Bei ν* liegt auf dem Randbogen \textbar τ\textbar{} = 1 eine Barriere von 7,89851 auf 7,90956 bei arg τ ≈ 72°, also 0,14 \%.
\item
  \textbf{Neu gegenüber dem Pilot: Ein zweiter, kontinuierlicher Übergang bei ν₂ = 8,606} \textbf{{[}nachgeprüft{]}}. Dort wird das Quadratgitter in Streckrichtung instabil, und ein Rechteckgitter τ = i·y übernimmt. Sein Seitenverhältnis wächst stetig: 1,056 (ν = 9), 1,100 (10), 1,141 (12), 1,170 (15), 1,194 (20), 1,215 (30). Die Pilotaussage „Quadrat für alle ν \textgreater{} 3,918`` ist damit \textbf{falsch} für ν \textgreater{} 8,606. Der Pilot hatte nur bis ν = 6 gerechnet.
\item
  \textbf{Grenzfall ν → ∞ {[}abgeleitet, nicht bewiesen{]}:} Für große ν zählt nur das kleinste Dreieck. Im Rechteckgitter konkurrieren das rechtwinklige Dreieck (Produkt √(y + 1/y)) und das kollineare Dreieck (Produkt 2y\textsuperscript{−3/2}). Gleichsetzen ergibt y⁴ + y² = 4, also \textbf{y∞ = √((√17 − 1)/2) ≈ 1,2496}. Die numerische Folge läuft darauf zu. Für ν ≥ 50 ist die FFT-Summe wegen des Dynamikumfangs nicht mehr genau genug; diese Werte sind verworfen.
\end{itemize}

\begin{figure}
\centering
\includegraphics{../verification/figures/fig_2d_energy_difference.pdf}
\caption{Energieunterschied Quadrat bzw. bestes Rechteck zum Hexagonalgitter. Die Nulldurchgänge markieren die beiden Übergänge.}
\end{figure}

\begin{figure}
\centering
\includegraphics{../verification/figures/fig_2d_stability.pdf}
\caption{Kleinster Hesse-Eigenwert von Hexagonal- und Quadratgitter. Grau: Bereich, in dem beide lokal stabil sind (Koexistenz, Übergang erster Ordnung). Rechts: Instabilität des Quadrats bei ν₂ = 8,606.}
\end{figure}

\begin{figure}
\centering
\includegraphics[width=1\textwidth,height=\textheight]{../verification/figures/fig_2d_landscape.pdf}
\caption{Energielandschaft log₁₀(T/T\_min − 1) über der Fundamentaldomäne bei vier Exponenten. Kreis: Hexagonalgitter, Quadrat: Quadratgitter. Bei ν = 10 liegt das Minimum auf der imaginären Achse oberhalb von τ = i (Rechteckgitter).}
\end{figure}

\begin{figure}
\centering
\includegraphics{../verification/figures/fig_2d_rect_ratio.pdf}
\caption{Seitenverhältnis des optimalen Gitters für ν \textgreater{} ν₂.}
\end{figure}

\subsection{3D: BCC ist das beste gefundene Gitter, FCC wird instabil (V3, V5)}

Die globale Suche lief über alle 3D-Bravais-Gitter mit Einheitsvolumen, also einen 5-dimensionalen Formraum. Pro ν gab es 16 Startpunkte (FCC, BCC, SC und 13 zufällige), jeweils mit Nelder-Mead und LLL-Reduktion.

\begin{longtable}[]{@{}lllll@{}}
\toprule\noalign{}
ν & FCC & BCC & SC & bestes gefundenes Gitter \\
\midrule\noalign{}
\endhead
\bottomrule\noalign{}
\endlastfoot
3,2 & 52,0983519937 & \textbf{52,0006824132} & 56,0330699444 & BCC \\
3,5 & 38,1993132080 & \textbf{38,0743775719} & 41,9368217767 & BCC \\
4,0 & 25,2265705426 & \textbf{25,0532095478} & 28,6610705747 & BCC \\
4,5 & 17,9623763656 & \textbf{17,7406073980} & 21,0980682627 & BCC \\
5,0 & 13,3808991194 & \textbf{13,1146038631} & 16,2165790058 & BCC \\
6,0 & 8,0433635309 & \textbf{7,7101880730} & 10,2962433924 & BCC \\
7,0 & 5,1456778058 & \textbf{4,7819587144} & 6,8781813217 & BCC \\
9,0 & 2,3132634600 & \textbf{1,9784839245} & 3,2722630842 & BCC \\
12,0 & 0,7708588009 & \textbf{0,5646023215} & 1,1328205539 & BCC \\
\end{longtable}

\begin{itemize}
\tightlist
\item
  \textbf{BCC ist bei allen neun Exponenten das beste gefundene Gitter} \textbf{{[}nachgeprüft, numerisch{]}}. Alle Startpunkte, die tief genug kamen, liefen nach BCC. Die Werte bei ν = 3,5 bis 9 stimmen mit GZL auf 10⁻¹⁰ bis 10⁻¹⁵ überein (V1).
\item
  \textbf{BCC ist ein lokales Minimum} gegenüber \emph{allen} volumenerhaltenden Verzerrungen: Alle fünf Hesse-Eigenwerte sind positiv, geprüft bei ν = 3,5; 4,5; 6; 9. Auf dem Bain-Pfad (tetragonale Verzerrung) liegt das Minimum bei c/a = 1, also genau bei BCC. Das betrifft direkt die offene Stabilitätsfrage aus \href{https://arxiv.org/abs/2504.07338}{arXiv:2504.07338}, dort allerdings für Lennard-Jones plus ATM und nicht für die reine Dreieckssumme.
\item
  \textbf{Neu: FCC verliert bei ν = 3,752 seine lokale Stabilität} \textbf{{[}nachgeprüft{]}}. Der kleinste Hesse-Eigenwert ist 2,02 bei ν = 3,5, 0,013 bei 3,75 und −5,97 bei 4,0. Oberhalb davon ist FCC ein Sattelpunkt.
\item
  \textbf{Einschränkung:} Eine Suche mit 16 Starts ist kein Beweis des globalen Minimums. Unterhalb von ν = 3,2 wurde nicht gerechnet.
\end{itemize}

\begin{figure}
\centering
\includegraphics{../verification/figures/fig_3d_lattices.pdf}
\caption{Abstand von FCC und SC zu BCC in Prozent.}
\end{figure}

\subsection{Modulare Graphfunktionen: Identitäten auf Maschinengenauigkeit (V4)}

Beide Identitäten wurden mit der Referenzsumme im Impulsraum geprüft, ohne GZL und an fünf τ-Werten:

\begin{itemize}
\tightlist
\item
  \textbf{C₂,₂,₁ = (2/5)E₅ + ζ(5)/30} gilt bis auf \textbf{10⁻¹⁵} relativ \textbf{{[}nachgeprüft{]}}. Der Pilot kam über die ε-Extrapolation in GZL nur auf 10⁻⁸.
\item
  \textbf{C₁,₁,₁ = E₃ + ζ(3)} gilt bis auf \textbf{10⁻¹⁴} \textbf{{[}nachgeprüft{]}}. Das gelingt aber nur, wenn man in der Extrapolation die Terme R\textsuperscript{−2} log R berücksichtigt. Mit reinen Potenzen bleibt ein systematischer Fehler von 10⁻⁶. Der Grund: Alle drei Kanten liegen bei ν = d, dadurch divergiert die Nachbarsumme logarithmisch.
\item
  Für C₂,₂,₂, wo alle ν \textgreater{} d sind, stimmen GZL und Referenz auf 10⁻¹⁵ überein. Das Wörterbuch „planare MGF = GZL-Zeta des dualen Graphen`` ist damit am Dreieck/Theta-Graph bestätigt.
\item
  \textbf{Folgerung für Idee 1:} Wegen T\textsubscript{2s}(Λ\_τ) = π\textsuperscript{3s} C\textsubscript{s,s,s}(τ) übersetzen sich die 2D-Ergebnisse direkt. Das Minimum von C\textsubscript{s,s,s} über der Fundamentaldomäne springt bei \textbf{s* = 1,95918} von e\textsuperscript{iπ/3} nach i und wird für \textbf{s \textgreater{} 4,303} zu einem Rechteckpunkt.
\end{itemize}

\subsection{Was sich am Bericht ändert}

\begin{itemize}
\tightlist
\item
  Idee 1 ist bestätigt und wird stärker. Das 2D-Phasendiagramm hat \textbf{drei} Phasen (hexagonal → Quadrat → Rechteck), mit Übergängen erster und zweiter Ordnung. Dazu kommt in 3D ein klares Bild: BCC ist das beste Gitter, FCC wird ab ν = 3,752 instabil.
\item
  Die Pilotaussage „Quadrat für ν \textgreater{} 3,918`` ist für ν \textgreater{} 8,606 korrigiert.
\item
  Die Präzisionsangabe für die MGF-Identitäten steigt von 10⁻⁸ auf 10⁻¹⁴ bis 10⁻¹⁵. Diese Präzision stammt aber aus der eigenen Referenzsumme und nicht aus GZL. Die Aussage in Idee 3, dass GZL für a = 1-Kanten eine Erweiterung braucht, bleibt richtig.
\end{itemize}

\subsection{Neuheitsprüfung}

Zwei unabhängige Such-Agenten haben die Mathematik- und die Physik/Chemie-Literatur gezielt nach diesen Ergebnissen durchsucht. Die Notizen stehen in \texttt{research\_notes/…/07\_novelty\_math.md} und \texttt{08\_novelty\_physics.md}. Eine Einschränkung vorweg: Semantic Scholar, die arXiv-API und Google Scholar waren nicht oder nur eingeschränkt erreichbar. Gesucht wurde deshalb über die arXiv-Suche, eine Websuche und Volltext-Durchsicht der Schlüsselarbeiten.

\textbf{Nächste Vorarbeit.} Robles-Navarro, Cooper, Buchheit, Busse, Burrows, Smits und Schwerdtfeger (\href{https://arxiv.org/abs/2504.07338}{arXiv:2504.07338}, J. Chem. Phys. 163, 094104, 2025) definieren genau diese Summe als „three-body zeta function`` ζ⁽³⁾(ν₁,ν₂,ν₃). Sie verwenden sie aber nur bei \textbf{ν = 3}, als Radialteil der ATM-Energie und zusammen mit Lennard-Jones. In 2D vergleichen sie nur Quadrat und Hexagon; in 3D untersuchen sie nur den Bain-Pfad und zeigen dort (Anhang H), dass BCC für beliebige Exponenten stationär ist. \textbf{Externe Kontrolle {[}nachgeprüft{]}:} Deren veröffentlichte ν = 3-Werte, auf Einheitsdichte umgerechnet (2D: hex 13,39179, Quadrat 13,65289; 3D: BCC ≈ 66,714, FCC ≈ 66,794), stimmen mit unserer Referenz auf alle angegebenen Stellen überein (13,3917941102; 13,6528937154; 66,7135509; 66,7943555; \texttt{results/v6\_external\_check\_nu3.log}).

\begin{longtable}[]{@{}
  >{\raggedright\arraybackslash}p{(\columnwidth - 4\tabcolsep) * \real{0.3996}}
  >{\raggedright\arraybackslash}p{(\columnwidth - 4\tabcolsep) * \real{0.3996}}
  >{\raggedright\arraybackslash}p{(\columnwidth - 4\tabcolsep) * \real{0.1708}}@{}}
\toprule\noalign{}
\begin{minipage}[b]{\linewidth}\raggedright
Ergebnis
\end{minipage} & \begin{minipage}[b]{\linewidth}\raggedright
Urteil der Neuheitsprüfung
\end{minipage} & \begin{minipage}[b]{\linewidth}\raggedright
Sicherheit
\end{minipage} \\
\midrule\noalign{}
\endhead
\bottomrule\noalign{}
\endlastfoot
2D: Sprung hexagonal → Quadrat bei ν* = 3,9184, erster Ordnung, Koexistenzbereich 3,806--4,278 & \textbf{keine Vorarbeit gefunden} (nur ν = 3-Daten mit anderer Normierung und Energie) & mittel--hoch (\textasciitilde85 \%) \\
2D: kontinuierlicher Übergang Quadrat → Rechteck bei ν₂ = 8,606, Seitenverhältnis → \textasciitilde1,25 & \textbf{keine Vorarbeit gefunden}. Verwandt: Übergänge Quadrat → Rechteck für \emph{Paar}potentiale (\href{https://arxiv.org/abs/2312.01395}{arXiv:2312.01395}) und eine Rechteckverzerrung durch \emph{anderen} Mechanismus in 2504.07338 & mittel--hoch (\textasciitilde85 \%) \\
Minimum von C\textsubscript{s,s,s}(τ): springt bei s* = 1,959, Rechteckpunkt für s \textgreater{} 4,303. Damit liegt das Minimum von C₂,₂,₂, C₃,₃,₃, C₄,₄,₄ bei τ = i und für s ≥ 5 bei τ = iy, y \textgreater{} 1 & \textbf{keine Vorarbeit gefunden} (Volltexte der MGF-Hauptarbeiten durchgesehen). Nur s = 1 ist klassisch (über C₁,₁,₁ = E₃ + ζ(3)) & mittel (\textasciitilde75 \%) \\
3D: BCC stationär und Minimum auf dem Bain-Pfad & \textbf{bekannt} (2504.07338, Anhang H) & hoch \\
3D: BCC bestes Gitter über \emph{alle} Bravais-Gitter für ν = 3,2--12; FCC ab ν = 3,752 instabil (Sattel) & \textbf{keine Vorarbeit gefunden} & mittel (\textasciitilde65 \%) \\
Planare MGF = Graph-Zeta des dualen Graphen & im Kern \textbf{per Konstruktion bekannt} (dieselbe Gittersumme, andere Normierung). Neu ist nur die Nutzung als Rechenwerkzeug & -- \\
\end{longtable}

\textbf{Verwandte, aber andere Arbeiten.}

\begin{itemize}
\tightlist
\item
  Für \textbf{Paar}energien gilt das Gegenteil: FCC ist optimal (\href{https://arxiv.org/abs/2609.17356}{arXiv:2609.17356}).
\item
  Paar-Analoga mit Phasenfolgen hexagonal → rhombisch → Quadrat → Rechteck gibt es von Luo und Wei. Unsere Dreikörperenergie springt dagegen direkt von hexagonal zu Quadrat.
\item
  Dreikörper-Kristallisation mit \textbf{Winkel}potentialen (Stefanelli, Friedrich, Kreutz u. a.) und Embedded-Atom-Modelle (Bétermin, Friedrich, Stefanelli, arXiv:2101.05602) behandeln andere Energien.
\item
  Gaußsche Dreikörper-Abstoßung in 2D (Sellin \& Babaev, arXiv:1308.2109) erzeugt Streifenphasen, kein Quadratgitter.
\end{itemize}

\textbf{Fazit.} Soweit sich das mit einer Literatursuche feststellen lässt, sind \textbf{der 2D-Phasenverlauf hexagonal → Quadrat → Rechteck für die reine Dreikörper-Potenzgesetzenergie, die Aussagen über die Minima von C\textsubscript{s,s,s} und das globale 3D-Ergebnis neu}. Die Neuheit muss man eng formulieren: reine Produktform, feste Dichte, alle Bravais-Gitter, allgemeines ν. Das größte Risiko ist nicht eine existierende Arbeit, sondern die Gruppe Buchheit/Schwerdtfeger. Zwei ihrer Arbeiten (\href{https://arxiv.org/abs/2504.11989}{arXiv:2504.11989}, \href{https://arxiv.org/abs/2609.18918}{arXiv:2609.18918}) kündigen eine „ongoing investigation into the influence of many-body interactions on the stability of crystal lattices`` an. Eine frühe Kontaktaufnahme ist deshalb der wichtigste nächste Schritt.

\textbf{Offene Punkte.}

\begin{enumerate}
\def\labelenumi{\arabic{enumi}.}
\tightlist
\item
  Die Summe konvergiert schon für ν \textgreater{} 2d/3, nicht erst für ν \textgreater{} d; die Grenze ν \textgreater{} d stammt aus der GZL-Implementierung. Bei ν = 2 folgt das hexagonale Optimum exakt aus C₁,₁,₁ = E₃ + ζ(3). Der Bereich 4/3 \textless{} ν \textless{} 2 ist nicht gerechnet, weil dort die Abschneidefehler der Referenz sehr langsam abfallen.
\item
  Die Klassiker zu Edelgaskristallen mit ATM (Jansen \& Lombardi 1966 u. a.) wurden nur über Metadaten geprüft.
\item
  Die 3D-Aussage beruht auf 16 Starts pro ν und ist kein Beweis.
\end{enumerate}

\begin{center}\rule{0.5\linewidth}{0.5pt}\end{center}

\section{Quellenliste}

\textbf{Software und Dokumentation}: \href{https://github.com/graph-zeta/gzl}{GZL GitHub (README)}, \href{https://github.com/graph-zeta/gzl/blob/main/DOCUMENTATION.md}{DOCUMENTATION.md}, \href{https://pypi.org/project/gzl/}{gzl auf PyPI}, \href{https://pypi.org/project/epsteinlib/}{epsteinlib auf PyPI}.

\textbf{Eigene Pilotrechnungen (nicht begutachtet)}: \texttt{research/pilots/}:

\begin{itemize}
\tightlist
\item
  \texttt{mgf\_test.py}: C₂,₂,₁- und Zagier-Test mit ε-Verschiebung;
\item
  \texttt{mgf\_test2.py}: Richardson-Extrapolation und tetraedrische K₄-MGF;
\item
  \texttt{opt\_test.py}: 2D-Gitterscan und 3D FCC/BCC/SC;
\item
  \texttt{opt2.py}: Nelder-Mead und Brute-Force-Check;
\item
  \texttt{opt3.py}: Brent-Wurzel ν* und Barriere entlang des Bogens.
\end{itemize}

\textbf{Nachprüfung (unabhängige Referenz, Abschnitt 7)}: \texttt{research/verification/}:

\begin{itemize}
\tightlist
\item
  \texttt{lattice\_sums.py}: Referenzimplementierung der Dreieckssumme ohne GZL-Code;
\item
  \texttt{v1\_gzl\_vs\_reference.py}: GZL gegen Referenz, 55 Fälle;
\item
  \texttt{v2\_2d\_landscape.py}, \texttt{v2b\_high\_nu.py}, \texttt{v2c\_large\_nu.py}: 2D-Phasendiagramm, Stabilität, zweiter Übergang;
\item
  \texttt{v3\_3d\_lattices.py}, \texttt{v5\_3d\_bain\_stability.py}, \texttt{v5b\_fcc\_threshold.py}: globale 3D-Suche, Bain-Pfad, Hesse-Matrizen;
\item
  \texttt{v4\_modular\_graph\_functions.py}, \texttt{v4b\_c111\_log\_fit.py}: MGF-Identitäten;
\item
  \texttt{make\_figures.py}: Abbildungen; Rohdaten und Logs in \texttt{results/}.
\end{itemize}

\begin{longtable}[]{@{}
  >{\raggedright\arraybackslash}p{(\columnwidth - 4\tabcolsep) * \real{0.1768}}
  >{\raggedright\arraybackslash}p{(\columnwidth - 4\tabcolsep) * \real{0.6916}}
  >{\raggedright\arraybackslash}p{(\columnwidth - 4\tabcolsep) * \real{0.1016}}@{}}
\toprule\noalign{}
\begin{minipage}[b]{\linewidth}\raggedright
arXiv
\end{minipage} & \begin{minipage}[b]{\linewidth}\raggedright
Kurzangabe (laut Notizen)
\end{minipage} & \begin{minipage}[b]{\linewidth}\raggedright
Jahr
\end{minipage} \\
\midrule\noalign{}
\endhead
\bottomrule\noalign{}
\endlastfoot
\href{https://arxiv.org/abs/2609.18918}{2609.18918} & Buchheit \& Rupp, Methode der Graph-Zetafunktionen & 2026 \\
\href{https://arxiv.org/abs/2609.18761}{2609.18761} & Duft, Adelhardt, Koziol, Buchheit, Schmidt, „Exact and fast series expansions for quantum models with long-range interactions'' & 2026 \\
\href{https://arxiv.org/abs/2609.28282}{2609.28282} & Buchheit \& Busse, anisotrope singuläre Summen aus Ableitungen der Epstein-Zeta & 2026 \\
\href{https://arxiv.org/abs/2504.11989}{2504.11989} & Buchheit \& Busse, „Epstein zeta method for many-body lattice sums'' & 2025 \\
\href{https://arxiv.org/abs/2412.16317}{2412.16317} & Buchheit, Busse, Gutendorf, EpsteinLib & 2024 \\
\href{https://arxiv.org/abs/2504.07338}{2504.07338} & Robles-Navarro et al., LJ + ATM auf kuboidalen Gittern & 2025 \\
\href{https://arxiv.org/abs/2012.05413}{2012.05413} & Analytische LJ-Schwingungseffekte, Einstein-Modell & 2020 \\
\href{https://arxiv.org/abs/2403.00421}{2403.00421} & Adelhardt, Koziol, Langheld, Schmidt, MC-Techniken für LR-Quantenmagnete & 2024 \\
\href{https://arxiv.org/abs/1802.06684}{1802.06684} & Fey, Kapfer, Schmidt, 2D-LR-Quantenmagnete & 2018 \\
\href{https://arxiv.org/abs/2203.08081}{2203.08081} & Langheld et al., Skalierung oberhalb der oberen kritischen Dimension & 2022 \\
\href{https://arxiv.org/abs/2408.13145}{2408.13145} & Adelhardt, Duft, Schmidt, LR-XXZ-Bilayer & 2024 \\
\href{https://arxiv.org/abs/2604.12754}{2604.12754} & Adelhardt, Muleady, Schmidt, Gorshkov, LR-Spin-1-Kette & 2026 \\
\href{https://arxiv.org/abs/2604.20831}{2604.20831} & LR-Spin-1-Kette, QMC α\textsubscript{c} = 2,49(1) & 2026 \\
\href{https://arxiv.org/abs/2601.20058}{2601.20058} & LR-XY-Kette, α\textsubscript{c} ≈ 2,7 & 2026 \\
\href{https://arxiv.org/abs/2606.22407}{2606.22407} & Li, Fan, Chen, Deng, Universalitätsdiagramm LR-Quantenkritikalität & 2026 \\
\href{https://arxiv.org/abs/2512.04805}{2512.04805} & Xiao, Liu, Fan, Deng, 2D LR-Ising σ* = 2 & 2025 \\
\href{https://arxiv.org/abs/2512.01956}{2512.01956} & Xiao et al., 2D LR-Heisenberg & 2025 \\
\href{https://arxiv.org/abs/2608.20750}{2608.20750} & Liu, Xiao, Fan, Deng, 2D LR-Perkolation & 2026 \\
\href{https://arxiv.org/abs/2508.18807}{2508.18807} & Hutchcroft, Critical LRP I & 2025 \\
\href{https://arxiv.org/abs/2508.18808}{2508.18808} & Hutchcroft, Critical LRP II & 2025 \\
\href{https://arxiv.org/abs/2404.07276}{2404.07276} & Hutchcroft, Dreiecksbedingung LRP & 2024 \\
\href{https://arxiv.org/abs/2607.22387}{2607.22387} & Blanc-Renaudie \& Hutchcroft, k-Punkt-Funktion & 2026 \\
\href{https://arxiv.org/abs/2602.07818}{2602.07818} & Zweischleifen-ε-Entwicklung LR-O(n) & 2026 \\
\href{https://arxiv.org/abs/2608.15120}{2608.15120} & 6 − ε für LR-Perkolation & 2026 \\
\href{https://arxiv.org/abs/2510.02458}{2510.02458} & Pagni et al., FRG für die 1D-LR-Ising-Kette & 2025 \\
\href{https://arxiv.org/abs/2604.07130}{2604.07130} & Okuyama \& Ohzeki, 1D-LR-Spinglas & 2026 \\
\href{https://arxiv.org/abs/2412.08814}{2412.08814} & Liu, Park, Slade, FSS oberhalb d\textsubscript{c} & 2024 \\
\href{https://arxiv.org/abs/2308.09709}{2308.09709} & NQS für die AF-LR-Ising-Kette & 2023 \\
\href{https://arxiv.org/abs/2609.17356}{2609.17356} & Luo \& Wei, Sarnak--Strömbergsson & 2026 \\
\href{https://arxiv.org/abs/2605.07580}{2605.07580} & Luo \& Wei, Minimierer von Theta- und Epstein-Verhältnissen & 2026 \\
\href{https://arxiv.org/abs/2411.17199}{2411.17199} & Deng \& Luo, hexagonale Optimalität & 2024 \\
\href{https://arxiv.org/abs/2407.20762}{2407.20762} & Bétermin \& Furlanetto, p-Normen & 2024 \\
\href{https://arxiv.org/abs/2312.01395}{2312.01395} & Bétermin, Šamaj, Travěnec, Strukturübergänge & 2023 \\
\href{https://arxiv.org/abs/1512.06779}{1512.06779} & MGFs und elliptische Polylogarithmen & 2015 \\
\href{https://arxiv.org/abs/1608.04393}{1608.04393} & D'Hoker \& Kaidi, MGF-Identitäten & 2016 \\
\href{https://arxiv.org/abs/2502.05531}{2502.05531} & Claasen \& Doroudiani, iterierte Eisenstein-Integrale & 2025 \\
\href{https://arxiv.org/abs/2511.15883}{2511.15883} & Schlotterer, Sohnle, Tao, elliptische MGFs & 2025 \\
\href{https://arxiv.org/abs/2512.21413}{2512.21413} & Fedosova \& Klinger-Logan, Konvolutionsidentitäten & 2025 \\
\href{https://arxiv.org/abs/2603.20550}{2603.20550} & Dobrowolski, Mordell--Tornheim & 2026 \\
\href{https://arxiv.org/abs/2607.25595}{2607.25595} & Portner, Feynman-Perioden und svMZVs & 2026 \\
\href{https://arxiv.org/abs/2603.20372}{2603.20372} & 256-Qubit-Simulation TmMgGaO₄ & 2026 \\
\href{https://arxiv.org/abs/2608.07178}{2608.07178} & Pasqal/Browaeys, 2D-TFIM mit 256 Qubits & 2026 \\
\href{https://arxiv.org/abs/2601.16275}{2601.16275} & Caltech, DSF einer Rydberg-Kette & 2026 \\
\href{https://arxiv.org/abs/2505.09884}{2505.09884} & NaTmSe₂ als TFIM & 2025 \\
\href{https://arxiv.org/abs/2306.03544}{2306.03544} & KTmSe₂-INS & 2023 \\
\href{https://arxiv.org/abs/2311.11726}{2311.11726} & Chen et al., Quench-Spektroskopie dipolar XY & 2023 \\
\href{https://arxiv.org/abs/2605.07685}{2605.07685} & DMRG dipolares XY auf Archimedischen Gittern & 2026 \\
\href{https://arxiv.org/abs/2605.28599}{2605.28599} & NLCE+QA auf Ionen-QPU & 2026 \\
\href{https://arxiv.org/abs/2512.09037}{2512.09037} & Magnon-Bindungszustände im 2D-LRTFIM & 2025 \\
\href{https://arxiv.org/abs/2605.13969}{2605.13969} & Squeezing-Übergang in Potenzgesetz-Bilayern & 2026 \\
\href{https://arxiv.org/abs/2605.27484}{2605.27484} & Koziol, ferromagnetisches Dicke-Ising-Modell & 2026 \\
\href{https://arxiv.org/abs/2604.22743}{2604.22743} & Fitzner, Lesanovsky, Sbierski, Thermometrie Kagome-Rydberg & 2026 \\
\href{https://arxiv.org/abs/2608.15576}{2608.15576} & Optimale Abschneideschranke O(R\textsuperscript{−(p−d)}) & 2026 \\
\href{https://arxiv.org/abs/2606.20522}{2606.20522} & iMPS ohne Surrogat für algebraische Schwänze & 2026 \\
\href{https://arxiv.org/abs/2512.24596}{2512.24596} & Xie \& Schotland, Photonbandstruktur atomarer Gitter & 2025 \\
\href{https://arxiv.org/abs/2609.17313}{2609.17313} & Exzeptionelle Topologie in 2D-Atomarrays & 2026 \\
\href{https://arxiv.org/abs/2507.18399}{2507.18399} & Modifizierte Lüscher-Zeta mit langreichweitiger Kraft & 2025 \\
\href{https://arxiv.org/abs/1702.01296}{1702.01296} & Matzelle \& Tiburzi, QED-Finite-Volume-Summen & 2017 \\
\href{https://arxiv.org/abs/2608.10041}{2608.10041} & He \& Hu, Madelung-Randterm & 2026 \\
\href{https://arxiv.org/abs/2503.00977}{2503.00977} & Calara \& Miller, Madelung per Multipolsummation & 2025 \\
\href{https://arxiv.org/abs/2603.13086}{2603.13086} & Erb, Meyer, Ottnad, Photonpropagator Gitter-QCD & 2026 \\
\href{https://arxiv.org/abs/2606.06617}{2606.06617} & PSWF-LR, langreichweitige ML-Potenziale & 2026 \\
\href{https://arxiv.org/abs/2512.18029}{2512.18029} & Latent Ewald & 2025 \\
\href{https://arxiv.org/abs/2308.03140}{2308.03140} & libMBD & 2023 \\
\href{https://arxiv.org/abs/2412.01076}{2412.01076} & LR-Kitaev-Ketten nach Quench & 2024 \\
\href{https://arxiv.org/abs/2609.28735}{2609.28735} & Erweiterte LR-Kitaev-Modelle & 2026 \\
\href{https://arxiv.org/abs/2606.13499}{2606.13499} & 2D-LR-XY mit 200 Ionen & 2026 \\
\end{longtable}

\end{document}
